\chapter{Satellite Communications}
\label{ch:satellite}

\begin{chapterintro}
A communication satellite is a repeater on a very tall mast: tall enough to see a third of the
planet, or to sweep past a user every few minutes. That single idea---first written down by
Arthur C.~Clarke in 1945---has carried transoceanic telephone calls, the world's television,
ships' and aircraft's links, emergency messages from phones in the wilderness, broadband to
millions of homes beyond the reach of fibre, and pictures from the edge of the solar system.
This chapter is the satellite engineer's toolkit. We start with the history, because every
generation of satellites was shaped by the lessons of the one before. We then develop the
orbital mechanics a communicator needs (periods, geometry, look angles, delay and Doppler),
the spectrum and its regulation, and the satellite link budget in depth, including rain
fades and interference. We open up the spacecraft---antennas, spot beams, transponders and
their nonlinear amplifiers---and study the multiple-access schemes and the DVB-S2 standard
that run through them. The last part of the chapter tours today's systems: VSAT networks,
high-throughput GEO satellites, low-Earth-orbit mega-constellations with laser crosslinks,
5G non-terrestrial networks and phones that talk to satellites directly, and finally the
Deep Space Network, where every decibel is fought for across billions of kilometres.
\end{chapterintro}

\section{Why satellites?}
\label{sec:ch23:why}

Every other chapter of this book fights distance with more infrastructure: more cable, more
towers, more cells. A satellite fights it with altitude. A single geostationary satellite sees
about 42\% of the Earth's surface; a single transmitter on it can therefore reach every home
in a continent at once, and the cost of serving one more receiver in that footprint is
precisely zero. That is the first and oldest advantage: \emph{broadcast}. The second is
\emph{reach}: a satellite serves ships in mid-ocean, aircraft at 11\,km, oil platforms,
Antarctic stations, disaster zones whose terrestrial network has been destroyed, and
villages that no fibre operator will ever reach profitably. The third is \emph{independence
from the ground}: a satellite link needs nothing between its two ends, which is why armies,
news crews and emergency services rely on it.

Satellites also have three fundamental handicaps, and much of this chapter is about living
with them. They are \emph{far away}, so the signal is weak---a geostationary downlink loses
about 205\,dB to spreading alone---and the round trip takes about half a second. They are
\emph{expensive to build, launch and impossible to repair}, so every kilogram and every watt
on board is precious and every component must survive fifteen years of radiation and thermal
cycling without maintenance. And they \emph{share a finite resource}---the spectrum and, for
geostationary satellites, a single ring of orbital positions above the equator---with every
other satellite in the sky, so their design is bound by international regulation to an
extent unknown in most of engineering.

The history of satellite communication is the history of engineers trading these advantages
and handicaps against each other. Low orbits cut the delay and the path loss but need many
satellites and tracking antennas; the geostationary orbit gives fixed antennas and wide
coverage at the cost of delay and path loss. Big satellites with big antennas let ground
terminals be small and cheap; small satellites need big ground antennas. The balance has
swung back and forth for sixty years, and it is swinging again as you read this.

\section{A short history of communication satellites}
\label{sec:ch23:history}

\begin{history}[Clarke's ``extra-terrestrial relays'']
In the October 1945 issue of \emph{Wireless World}, Arthur C.~Clarke---then an RAF radar
officer, later famous as a science-fiction writer---published a short article titled
``Extra-Terrestrial Relays''. He pointed out that a body orbiting 35\,800\,km above the equator
would circle the Earth once a day and so hang motionless in the sky; that three such stations
spaced 120$^\circ$ apart could provide worldwide coverage for radio and television; and that
the rockets being developed in wartime Germany would eventually make it possible to put them
there. Clarke imagined crewed stations with vacuum-tube equipment and solar power; the
orbit he described now bears his name, the \emph{Clarke orbit}. He never patented the idea,
and later remarked with good humour that he had given away a fortune. The engineering
arithmetic in the article---power, path loss and coverage---is recognisably the link budget of
Section~\ref{sec:ch23:linkbudget}.
\end{history}

\subsection{From Sputnik to Telstar}

The Space Age began on 4 October 1957, when Sputnik~1 transmitted its beacon on about 20 and
40\,MHz. It carried no communication payload, but it proved that orbit was reachable---and,
as Chapter~\ref{ch:history} recounts, its Doppler shift inspired satellite navigation. The
first voice relayed through space came a year later: Project SCORE, launched in December
1958, broadcast a tape-recorded Christmas message from President Eisenhower. Its batteries
lasted about two weeks.

Two design philosophies were then tested side by side. \textbf{Passive satellites}\index{passive satellite} simply
reflect signals: NASA's Echo~1 (1960) was a 30\,m aluminised plastic balloon at about 1600\,km.
Echo worked---telephone calls bounced off it between New Jersey and California---but the path
loss of a double spreading (up to the balloon and back from a sphere that reradiates in all
directions) made it hopeless for commercial traffic. The horn antenna that Bell Labs built at
Holmdel, New Jersey to receive Echo's feeble signals was so quiet that in 1964 Arno Penzias
and Robert Wilson, measuring its noise temperature, found about 3.5\,K of unexplained excess
noise coming from every direction: the cosmic microwave background, the afterglow of the Big
Bang. It is a satisfying fact for a communications engineer that one of the great discoveries
of physics began as a noise-temperature budget (Chapter~\ref{ch:noise}).

\textbf{Active satellites}\index{active satellite} receive, amplify and retransmit. AT\&T's Telstar~1, launched on
10 July 1962, was the first active commercial communication satellite. It received at about
6\,GHz and transmitted at about 4\,GHz---the ``6/4\,GHz'' C-band pairing still used today---through a
travelling-wave-tube amplifier, the ancestor of the tubes in Section~\ref{sec:ch23:payload}.
Within two weeks it carried the first live transatlantic television between Andover, Maine and
Pleumeur-Bodou in Brittany. But Telstar was in an elliptical low orbit (roughly 950 by
5900\,km): each ground station could use it for only about 20 minutes per pass, with huge
tracking antennas, and it failed in November 1962 after radiation damage to its transistors,
aggravated by the Starfish Prime high-altitude nuclear test of July 1962 which had pumped up
the Van Allen belts.

\subsection{The geostationary era}

The engineer who made Clarke's orbit practical was Harold Rosen of Hughes Aircraft. His team's
Syncom was small enough (a few tens of kilograms in orbit) to be launched to 36\,000\,km by the rockets of
the day, and it was \emph{spin-stabilised}: the whole satellite rotated like a gyroscope to hold
its orientation, avoiding the need for a complex attitude-control system. Syncom~2 (July
1963) reached a geosynchronous but inclined orbit; Syncom~3 (August 1964) was the first
truly geostationary satellite and relayed the Tokyo Olympic Games to the United States.

The Communications Satellite Act of 1962 had created COMSAT in the United States, and in 1964
an international consortium, \term{Intelsat}, was formed to own a global system. Its first
satellite, Intelsat~I (``Early Bird''), began commercial service across the Atlantic in 1965
with 240 telephone circuits \emph{or} one television channel. Each new generation multiplied
capacity: Intelsat~III (1968--70) completed global coverage over three ocean regions;
Intelsat~IV (1971) carried thousands of circuits in twelve 36\,MHz transponders---the
bandwidth that became the industry's unit of capacity. Satellites carried most intercontinental
telephone traffic until the optical submarine cables of the late 1980s (Chapter~\ref{ch:wireline})
took it back with a fraction of the delay.

The Soviet Union, whose territory lies mostly at high latitudes where geostationary
satellites sit low on the horizon, chose a different orbit. The \term{Molniya} (``lightning'')
satellites, first launched in 1965, used a highly elliptical 12-hour orbit with an apogee over
the northern hemisphere, where the satellite lingers for eight hours or so of each orbit
(Section~\ref{sec:ch23:orbits}). From 1967 the Orbita network of ground stations distributed
Moscow television across eleven time zones through them.

Domestic systems followed: Canada's Anik~A1 (1972) was the first domestic commercial
geostationary satellite; the United States had Westar and Satcom by the mid-1970s. In 1975 the
cable channel HBO began distributing its programme to cable head-ends by satellite, and within
a few years every cable network did the same---the start of the business that still pays for
most of the geostationary fleet. NASA's ATS-6 (1974) broadcast educational television directly
to thousands of simple village receivers in India in the SITE experiment of 1975--76, an early
demonstration of direct broadcasting.

\subsection{Television to the home, and mobility}

Direct-to-home (DTH) television needed two things: higher-power Ku-band satellites, so that
receiving dishes could shrink from several metres to well under one, and cheap low-noise
downconverters. In Europe, SES's Astra~1A (1988) carried Sky Television from 1989 to dishes
of about 60\,cm. In the United States, DirecTV launched in 1994 with an all-digital system
using MPEG-2 video compression and a concatenated Reed--Solomon and convolutional code; the same
year the European DVB project standardised \term{DVB-S}, which became the world standard for
digital satellite TV and was succeeded in 2005 by DVB-S2 (Section~\ref{sec:ch23:dvbs2}).
Digital compression multiplied the number of channels per transponder by five to ten, and
satellite TV became a mass market with hundreds of millions of subscribers.

Maritime communication followed a parallel path. Inmarsat, an intergovernmental
organisation founded in 1979, began service in 1982, replacing Morse and HF radio at sea
with L-band telephone and telex; later generations added aeronautical and land-mobile service,
the Global Maritime Distress and Safety System (GMDSS), and broadband (BGAN, Global Xpress).

\subsection{The first LEO boom and bust}

In the 1990s engineers returned to low orbits to serve \emph{handheld} terminals, whose
tiny antennas and batteries cannot close a link to 36\,000\,km. Motorola's \term{Iridium}
placed 66 active satellites in six near-polar planes at 780\,km, linked to each other by
Ka-band intersatellite links so that a call could be routed around the world in space.
Commercial service began in November 1998. Nine months later Iridium filed for bankruptcy,
with only tens of thousands of subscribers, having spent roughly five billion dollars; the system
was bought in 2000 by private investors for about 25 million dollars. The engineering was
superb; the market had changed underneath it, because terrestrial GSM had spread far faster
than anyone expected and a brick-sized, expensive satellite phone could not compete in the
cities where most calls were made. Globalstar (48 satellites at 1414\,km, bent-pipe,
CDMA based on IS-95) and the ``little LEO'' messaging system ORBCOMM also went through
bankruptcy. The most ambitious plan of the era, Teledesic's broadband constellation of first
840 and then 288 satellites, backed by Bill Gates and Craig McCaw, was abandoned in 2002
without a single operational satellite.

Iridium survived and prospered in niches---maritime, aviation, government, the Arctic---and a
second-generation Iridium NEXT constellation was launched in 2017--19. The lessons of the bust
were not that LEO was wrong, but that satellites had to be much cheaper, launch had to be much
cheaper, and the terminals had to be cheaper still.

\subsection{High-throughput satellites and the second LEO wave}

In the 2000s geostationary broadband was reinvented around \term{spot beams}: instead of one
wide beam covering a continent, the satellite forms dozens to hundreds of narrow Ka-band beams
and reuses the spectrum in each (Section~\ref{sec:ch23:antennas}). Anik~F2 (2004) and above all
ViaSat-1 (2011), which offered on the order of 140\,Gb/s---more than all previous North American
satellites combined---launched the era of \term{high-throughput satellites} (HTS). The ViaSat-3
generation aims at the order of a terabit per second per satellite. O3b (``the other three
billion''), from 2013, placed Ka-band satellites in a medium equatorial orbit at 8062\,km,
cutting round-trip delay to about 150\,ms for cellular backhaul and island connectivity.

Then launch costs collapsed with reusable rockets, and satellites began to be built on assembly
lines. SpaceX launched its first sixty Starlink satellites in May 2019 and by 2025 was operating
more than seven thousand of them, most in shells around 550\,km altitude, linked by laser
intersatellite links and serving millions of subscribers with flat electronically steered
user terminals. OneWeb (about 650 satellites at 1200\,km, bent-pipe) went bankrupt in 2020, was
rescued by the UK government and Bharti, and merged with Eutelsat in 2023. Amazon's Project
Kuiper, authorised for 3236 satellites, began deploying in 2025, and Chinese constellations of
similar scale are under way. Low Earth orbit, sixty years after Telstar, has become the busiest
neighbourhood in space.

The newest turn is \term{direct-to-device} service: satellites that talk to ordinary phones.
Apple's Emergency SOS via Globalstar (2022) sends short messages from an iPhone; AST
SpaceMobile's BlueWalker~3 test satellite (2022), with a phased array of about 64\,m$^2$, made
broadband calls to unmodified phones; SpaceX's direct-to-cell satellites carry LTE base stations
and began serving unmodified phones through a mobile operator's spectrum in 2024--25. 3GPP has
standardised \term{non-terrestrial networks} (NTN) since Release~17
(Section~\ref{sec:ch23:ntn}), so that satellites become just another kind of cell.

\begin{table}[htbp]\centering\small
\caption{Milestones of satellite communication.}\label{tab:ch23:milestones}
\begin{tabularx}{\linewidth}{@{}l l X@{}}
\toprule
Year & System & Significance\\
\midrule
1945 & Clarke, \emph{Wireless World} & Proposes geostationary relays for global coverage\\
1957 & Sputnik 1 & First artificial satellite; Doppler tracking leads to Transit and GPS\\
1958 & SCORE & First voice (recorded) relayed from orbit\\
1960 & Echo 1 & Passive reflector balloon; transcontinental calls\\
1962 & Telstar 1 & First active commercial satellite; live transatlantic TV\\
1964 & Syncom 3 & First geostationary satellite; Tokyo Olympics to the US\\
1965 & Intelsat I, Molniya 1 & Commercial GEO service; Soviet HEO system\\
1972--76 & Anik A1, Westar, ATS-6 & Domestic satellites; cable-TV distribution; direct broadcast trials\\
1982 & Inmarsat & Global maritime mobile service\\
1988--94 & Astra, DirecTV, DVB-S & Direct-to-home TV; digital compression\\
1998--2000 & Iridium, Globalstar & Handheld LEO voice; bankruptcies\\
2005 & DVB-S2 & LDPC codes, APSK, adaptive coding and modulation\\
2011 & ViaSat-1 & High-throughput Ka-band spot-beam satellite\\
2013 & O3b & MEO broadband for backhaul\\
2019-- & Starlink, OneWeb, Kuiper & LEO mega-constellations with phased-array terminals\\
2022-- & 3GPP Rel-17 NTN, D2D & Satellites as 5G cells; direct-to-phone service\\
2023 & DSOC on Psyche & Deep-space optical communication demonstrated\\
\bottomrule
\end{tabularx}
\end{table}

\section{Orbital mechanics for communicators}
\label{sec:ch23:orbits}

A communications engineer does not need to navigate a spacecraft, but must know where it is,
how fast it moves, how long it is visible and how its geometry changes the link. Fortunately
the essentials follow from one law of physics and a little trigonometry.

\subsection{Kepler and Newton}

Johannes Kepler distilled Tycho Brahe's planetary observations into three laws (1609--1619):
orbits are ellipses with the central body at one focus; the line from the central body to the
orbiter sweeps equal areas in equal times; and the square of the period is proportional to the
cube of the orbit's size. Newton showed that all three follow from an inverse-square
gravitational force. For a circular orbit of radius $r$ the derivation takes one line: gravity
supplies the centripetal force,
\begin{equation}
\frac{GMm}{r^2}=\frac{mv^2}{r}\quad\Longrightarrow\quad
v=\sqrt{\frac{\mu}{r}},\qquad
T=\frac{2\pi r}{v}=2\pi\sqrt{\frac{r^3}{\mu}},
\label{eq:ch23:period}
\end{equation}
where $\mu=GM_\oplus=3.986\times10^{14}$\,m$^3$/s$^2$ is the Earth's gravitational parameter
(known far more precisely than $G$ or $M_\oplus$ separately). For an ellipse the same formula
holds with $r$ replaced by the semi-major axis $a$---that is Kepler's third law---and the speed
at distance $r$ is given by the \emph{vis-viva} equation $v^2=\mu(2/r-1/a)$: fastest at perigee,
slowest at apogee.

\paragraph{The geostationary radius.} To stay above one point of the rotating Earth, a satellite
must orbit in the equatorial plane, eastward, with a period of one \emph{sidereal} day (the
Earth's rotation period relative to the stars), $T=86\,164.1$\,s, not the 86\,400\,s solar day.
Solving~\eqref{eq:ch23:period} for $r$,
\begin{equation}
r_{\text{GEO}}=\left(\frac{\mu T^2}{4\pi^2}\right)^{1/3}=42\,164\ \text{km},
\end{equation}
an altitude of 35\,786\,km above the equator, where the orbital speed is 3.07\,km/s.
Figure~\ref{fig:ch23_orbits_scale} shows how period, speed and delay vary with altitude and
marks some well-known systems. Periods range from about 90 minutes just above the atmosphere to
a day at GEO; speed falls from 7.7 to 3.1\,km/s; and the round-trip delay through a bent-pipe
satellite grows almost linearly with altitude, from a few milliseconds in LEO to half a second
at GEO.

\mdcfig{ch23_orbits_scale}{Orbital period, speed and bent-pipe round-trip time (user to
satellite to gateway and back, four slant ranges) versus altitude for circular orbits. Dots
mark the ISS (420\,km), Starlink ($\approx$550\,km), Iridium (780\,km), OneWeb (1200\,km),
O3b (8062\,km), GPS (20\,200\,km) and GEO (35\,786\,km). Computed with
\texttt{commlib.satellite}.}

\subsection{Orbital elements}

Six numbers, the \term{Keplerian orbital elements}, fix an orbit and the satellite's place on
it: the semi-major axis $a$ (size), the eccentricity $e$ (shape), the inclination $i$ (tilt of
the orbital plane to the equator), the right ascension of the ascending node $\Omega$ (where the
orbit crosses the equator going north, measured from a fixed direction in space), the argument
of perigee $\omega$ (where in the plane the perigee lies) and the mean anomaly $M$ at an epoch
(where the satellite is). The mean anomaly increases uniformly, $M=M_0+nt$ with mean motion
$n=\sqrt{\mu/a^3}$; the position on the ellipse follows from the eccentric anomaly $E$, which
solves \term{Kepler's equation}
\begin{equation}
M=E-e\sin E ,
\end{equation}
a transcendental equation that Newton's method solves in a few iterations. The true anomaly
$\nu$ and radius then follow from $\tan(\nu/2)=\sqrt{(1+e)/(1-e)}\tan(E/2)$ and
$r=a(1-e\cos E)$. Operators publish orbits in exactly this form: the \emph{two-line element}
sets (TLEs) distributed for tracked objects encode the six elements plus drag terms, and the
SGP4 propagator turns them into positions. The function \texttt{orbit\_eci} in
\texttt{commlib/satellite.py} implements the two-body version in fifteen lines.

\subsection{The orbit families}

\begin{table}[htbp]\centering\small
\caption{Orbit families used for communication (circular-orbit values unless stated).}
\label{tab:ch23:orbits}
\begin{tabularx}{\linewidth}{@{}l l l l X@{}}
\toprule
Orbit & Altitude & Period & RTT$^{*}$ & Examples and remarks\\
\midrule
LEO & 300--2000\,km & 90--127\,min & 4--26\,ms & Iridium, Globalstar, Starlink, OneWeb, Kuiper; many satellites, tracking terminals, handovers\\
MEO & 2000--35\,786\,km & 2--24\,h & 50--300\,ms & O3b (8062\,km), GPS, Galileo, GLONASS; the Van Allen belts are harsh on electronics\\
GEO & 35\,786\,km & 23\,h\,56\,min & 480--540\,ms & Broadcast, VSAT, HTS; fixed antennas; one crowded ring\\
Molniya & 500 $\times$ 39\,800\,km & 12\,h & varies & $e\approx0.74$, $i=63.4^\circ$; long dwell over high latitudes\\
Tundra & $e\approx0.25$--0.4 & 24\,h & varies & Geosynchronous but inclined and elliptical; used by Sirius satellite radio\\
\bottomrule
\multicolumn{5}{@{}l}{\footnotesize $^{*}$Bent-pipe user--satellite--gateway round trip, propagation only, from zenith to low elevation.}
\end{tabularx}
\end{table}

Table~\ref{tab:ch23:orbits} lists the families. \textbf{Low Earth orbit}\index{low Earth orbit} (LEO) lies between
the drag of the upper atmosphere, which below about 400\,km pulls satellites down within months
to a few years, and the inner Van Allen radiation belt. \textbf{Medium Earth orbit}\index{medium Earth orbit} (MEO) is home
to the navigation constellations and to O3b. \textbf{Geostationary orbit}\index{geostationary orbit} (GEO) is a single
circle. \textbf{Highly elliptical orbits}\index{highly elliptical orbit} (HEO) trade a fixed position for coverage of high
latitudes. Figure~\ref{fig:ch23_ground_tracks} shows the \term{ground tracks}---the paths of the
sub-satellite point---of three of them.

\mdcfig{ch23_ground_tracks}{Ground tracks. A 550\,km, 53$^\circ$-inclined LEO satellite
(three orbits) shifts about 24$^\circ$ westward each orbit as the Earth turns beneath it. The
12-hour Molniya orbit spends most of its time near apogee over high northern latitudes
(dots are one hour apart). A geosynchronous satellite with 8$^\circ$ inclination traces a
figure of eight; a geostationary one is a fixed point.}

\paragraph{Why 63.4$^\circ$?} The Earth is not a sphere: its equatorial bulge, described by
the coefficient $J_2\approx1.083\times10^{-3}$, adds a torque that makes orbits precess. Two
effects matter. The orbital plane rotates about the Earth's axis (\term{nodal regression}) at
\begin{equation}
\dot\Omega=-\tfrac32\,nJ_2\left(\frac{R_\oplus}{a(1-e^2)}\right)^2\cos i ,
\end{equation}
and the perigee rotates within the plane at
$\dot\omega=\tfrac34 nJ_2\,(R_\oplus/a(1-e^2))^2\,(5\cos^2 i-1)$. A Molniya orbit must keep its
apogee over the northern hemisphere, so $\dot\omega$ must vanish: $5\cos^2i=1$, or
$i=63.4^\circ$, the \emph{critical inclination}. The same $J_2$ effect is exploited for
\emph{sun-synchronous} orbits (inclinations of about 97--99$^\circ$ in LEO), whose plane turns
by 360$^\circ$ a year and so keeps a fixed angle to the Sun---the favourite of Earth-observation
satellites---and it is why LEO constellation shells at different inclinations slowly drift
relative to one another.

\subsection{Staying in place: perturbations, station keeping, eclipses}

A geostationary satellite is pushed off station by the Sun and Moon, which tilt its orbit by
roughly 0.75--0.95$^\circ$ per year, and by the Earth's slightly elliptical equator, which drags
it in longitude towards one of two stable points (near 75$^\circ$E and 105$^\circ$W). Operators
keep each satellite inside a box typically $\pm0.05^\circ$ to $\pm0.1^\circ$ on a side with
regular thruster burns; \emph{north--south} station keeping to cancel the inclination costs
most of the propellant, around 50\,m/s of velocity change per year. When the propellant runs
low, operators often stop north--south control and let the inclination grow: the satellite
then traces a daily figure of eight (Figure~\ref{fig:ch23_ground_tracks}) that tracking or
wide-beam ground antennas can tolerate, extending a revenue-earning life by years. Many modern
satellites use electric (ion or Hall-effect) thrusters whose high exhaust velocity cuts
propellant mass dramatically, at the cost of months of slow orbit raising after launch.

Twice a year, around the equinoxes, a GEO satellite passes through the Earth's shadow once per
day for up to about 70 minutes (the function \texttt{eclipse\_fraction} gives 69 minutes for
the umbra alone), and it must run from batteries. Around the same seasons a ground antenna
pointed at the satellite also sees the Sun directly behind it for a few minutes a day on
several consecutive days: this \term{sun outage} raises the antenna noise temperature by
thousands of kelvin and interrupts reception, and operators schedule around it. LEO satellites
are eclipsed on most orbits---up to about 36 minutes of each 96-minute orbit at 550\,km---so
battery cycling is a design driver.

\begin{pitfall}[Space is getting crowded]
There are now more than ten thousand active satellites, most in LEO, together with tens of
thousands of tracked pieces of debris and far more too small to track. A collision creates
fragments that can cause further collisions---the cascade described by Donald Kessler and
Burton Cour-Palais in 1978. Responsible operators keep LEO satellites low enough that drag
removes failed ones within a few years (the US FCC adopted a five-year post-mission disposal
rule in 2022), fit propulsion for active collision avoidance, and raise dying GEO satellites to
a graveyard orbit a few hundred kilometres above the geostationary ring. Mega-constellation
operators perform many thousands of avoidance manoeuvres a year. Designing for orderly disposal
is now part of a communication system's engineering, not an afterthought.
\end{pitfall}

\section{Geometry: seeing the satellite}
\label{sec:ch23:geometry}

\subsection{Elevation, slant range and coverage}

\begin{figure}[htbp]\centering
\begin{tikzpicture}[scale=1.0,>={Stealth[length=2.2mm]}]
\def\Re{2.6}
\coordinate (O) at (0,0);
\draw[navy,thick,fill=navy!5] (O) circle (\Re);
\coordinate (E) at (60:\Re);
\coordinate (S) at (32:6.2);
\draw[dashed,navy!60] (E) -- ($(O)!1.55!(E)$) node[above,font=\scriptsize,navy]{local vertical};
\draw[thick,navy] (O) -- node[left,font=\small,pos=0.55]{$R_\oplus$} (E);
\draw[thick,navy] (O) -- node[below right,font=\small,pos=0.62]{$r=R_\oplus+h$} (S);
\draw[very thick,accent] (E) -- node[above,font=\small,pos=0.55,sloped]{slant range $d$} (S);
\draw[gray,thick] ($(E)+(150:1.6)$) node[above,font=\scriptsize,gray]{local horizon} -- ($(E)+(-30:1.9)$);
\fill[navy] (O) circle (1.5pt) node[below left,font=\small]{$O$};
\fill[navy] (E) circle (1.5pt) node[above left,font=\small]{$E$};
\fill[accent] (S) circle (2.2pt) node[right,font=\small]{$S$ (satellite)};
\coordinate (P) at (32:\Re);
\fill[navy] (P) circle (1.2pt) node[below right,font=\scriptsize]{sub-satellite point};
\draw[thin] ($(O)+(32:0.9)$) arc (32:60:0.9); \node[font=\small] at (46:1.2) {$\gamma$};
\draw[thin] ($(E)+(-30:0.8)$) arc (-30:14.64:0.8);
\node[font=\small] at ($(E)+(-8:1.05)$) {$\theta$};
\draw[thin] ($(S)+(194.64:1.1)$) arc (194.64:212:1.1);
\node[font=\small] at ($(S)+(203:1.4)$) {$\eta$};
\end{tikzpicture}
\caption{Earth-station geometry. The Earth centre $O$, the station $E$ and the satellite $S$ form
a triangle with sides $R_\oplus$, $r$ and $d$. The elevation angle $\theta$ is measured from the
local horizon, $\gamma$ is the Earth central angle to the sub-satellite point and $\eta$ is the
off-nadir angle at the satellite; $\theta+\gamma+\eta=90^\circ$.}
\label{fig:ch23_geometry_tikz}
\end{figure}

Everything about the link depends on where the satellite appears in the sky. Consider the
triangle of Figure~\ref{fig:ch23_geometry_tikz}. The angle at the earth station $E$ between
the direction to the Earth's centre and the direction to the satellite is $90^\circ+\theta$,
where $\theta$ is the \term{elevation angle}. The law of cosines gives
$r^2=R_\oplus^2+d^2+2R_\oplus d\sin\theta$, a quadratic in the slant range $d$ whose positive
root is
\begin{equation}
d=\sqrt{r^2-R_\oplus^2\cos^2\theta}-R_\oplus\sin\theta .
\label{eq:ch23:slant}
\end{equation}
At zenith ($\theta=90^\circ$) this is simply the altitude $h$; at the horizon it is
$\sqrt{r^2-R_\oplus^2}$. The law of sines gives the \term{nadir angle} $\eta$ at which the
satellite must point its antenna, and the central angle $\gamma$ follows because the angles of
the triangle sum to 180$^\circ$:
\begin{equation}
\sin\eta=\frac{R_\oplus\cos\theta}{r},\qquad
\gamma=90^\circ-\theta-\eta=\arccos\!\left(\frac{R_\oplus}{r}\cos\theta\right)-\theta .
\label{eq:ch23:central}
\end{equation}

The area that sees the satellite above a minimum elevation $\theta_{\min}$ is a spherical cap
of angular radius $\gamma_{\max}$, whose area is a fraction $(1-\cos\gamma_{\max})/2$ of the
Earth's surface. For GEO with $\theta_{\min}=0$, $\gamma_{\max}=81.3^\circ$ and the fraction is
42.4\%; with a practical $\theta_{\min}=10^\circ$ it is 34\%, so three satellites cover all
but the polar caps, as Clarke said. The satellite sees the whole Earth within a nadir angle of
only $\pm8.7^\circ$, so a global beam is about 17$^\circ$ wide. For a 550\,km satellite and
$\theta_{\min}=25^\circ$ (a typical Starlink terminal constraint) the coverage radius
$R_\oplus\gamma_{\max}$ is only about 940\,km, a cap of about 0.54\% of the Earth, and the
satellite must steer its beams as far as 57$^\circ$ off nadir. Hundreds of satellites are
therefore needed just to have one in view everywhere, and thousands to provide capacity.
Figure~\ref{fig:ch23_geometry_curves} plots slant range, delay and coverage for representative
orbits. Note how strongly low elevation penalises LEO: at 550\,km the slant range doubles from
550\,km at zenith to 1120\,km at 25$^\circ$ (6\,dB more path loss) and exceeds 2700\,km at the
horizon.

\mdcfig{ch23_geometry_curves}{Left: slant range and one-way delay versus elevation for two
LEO altitudes, O3b's MEO and GEO. Right: fraction of the Earth's surface covered by one
satellite versus altitude, for three minimum elevation angles.}

\subsection{Look angles to a geostationary satellite}

For a GEO satellite at longitude $\lambda_s$ and an earth station at latitude $\phi$ and
longitude $\lambda$, let $\Delta\lambda=\lambda_s-\lambda$. The sub-satellite point is on the
equator, so spherical trigonometry gives the central angle directly, and from it the range,
elevation and azimuth (measured clockwise from true north):
\begin{gather}
\cos\gamma=\cos\phi\cos\Delta\lambda,\qquad
d=\sqrt{r^2+R_\oplus^2-2rR_\oplus\cos\gamma},\notag\\
\tan\theta=\frac{\cos\gamma-R_\oplus/r}{\sin\gamma},\qquad
\text{Az}=\operatorname{atan2}\!\left(\sin\Delta\lambda,\;-\sin\phi\cos\Delta\lambda\right).
\label{eq:ch23:lookangles}
\end{gather}
Figure~\ref{fig:ch23_geo_visibility} maps the elevation over the visible hemisphere. A station
on the satellite's meridian sees it at an elevation that falls almost linearly with latitude,
reaching zero at 81.3$^\circ$; above about 70$^\circ$ latitude GEO service becomes marginal,
which is why the Soviet Union chose Molniya and why Arctic users rely on Iridium, polar LEO
shells or HEO systems.

\mdcfig{ch23_geo_visibility}{Left: elevation angle to a geostationary satellite as a function
of the station's latitude and its longitude relative to the satellite. Right: elevation versus
latitude on the satellite's meridian and 30$^\circ$ and 60$^\circ$ away from it.}

\begin{worked}[Pointing a dish in Paris at Astra 19.2$^\circ$E]
A viewer in Paris ($\phi=48.86^\circ$N, $\lambda=2.35^\circ$E) wants to point a dish at the
Astra satellites at 19.2$^\circ$E, so $\Delta\lambda=16.85^\circ$.
Then $\cos\gamma=\cos48.86^\circ\cos16.85^\circ=0.658\times0.957=0.630$, so $\gamma=50.97^\circ$.
With $r=42\,164$\,km and $R_\oplus=6378$\,km,
$d=\sqrt{42\,164^2+6378^2-2(42\,164)(6378)(0.630)}\approx38\,470$\,km, a one-way delay of
128\,ms. The elevation is $\arctan[(0.630-0.151)/0.777]=31.6^\circ$ and the azimuth is
$\operatorname{atan2}(0.290,\,-0.753\times0.957)=158.1^\circ$: a little east of due south, about
a third of the way up the sky. The installer must also rotate the LNB to match the satellite's
linear polarisation as seen from Paris (the \emph{polarisation skew}), and remember that an
offset-fed dish looks physically tilted down by its offset angle. The call
\texttt{geo\_look\_angles(48.86, 2.35, 19.2)} returns these numbers.
\end{worked}

\subsection{Delay, and what it does to protocols}

Light takes 119\,ms to climb 35\,786\,km, and a typical GEO station sees its satellite at
38\,000--40\,000\,km, so a one-way hop through a bent-pipe satellite (up and down) takes
240--280\,ms and a request--response exchange, which needs two hops, 480--540\,ms before any
processing or queuing. The 3GPP study of non-terrestrial networks, TR~38.821, uses a worst-case
round trip of about 541\,ms for a transparent GEO satellite with both links at 10$^\circ$
elevation---exactly what~\eqref{eq:ch23:slant} gives. For telephony, ITU-T G.114 recommends
keeping one-way mouth-to-ear delay below about 150\,ms for good interactivity, with 400\,ms as a
planning limit; a single GEO hop fits, a double hop is painful, and every satellite voice
circuit needs excellent echo cancellation (Chapter~\ref{ch:pcm}). For data, delay hurts in two
ways.

\paragraph{Window-limited throughput.} A sender using a sliding-window protocol such as TCP
can have at most one window $W$ of unacknowledged data in flight, so its rate is limited to
\begin{equation}
R\le\frac{8W}{\text{RTT}}\ \text{b/s}.
\label{eq:ch23:tcpwindow}
\end{equation}
With the classic 64\,KiB TCP window and a 600\,ms round trip, $R\le0.87$\,Mb/s, no matter how
fast the satellite is. Window scaling (RFC~7323) removes this limit if both ends use large
buffers---the \emph{bandwidth--delay product} of a 100\,Mb/s GEO link is 7.5\,MB.

\paragraph{Slow start and handshakes.} Even with big windows, TCP probes the path by doubling
its congestion window every round trip, starting from about ten segments. Fetching a small
web object takes a handshake, a request and a handful of slow-start rounds, and each costs an
RTT. Figure~\ref{fig:ch23_tcp_delay} shows the consequence: a 1\,MB object needs about nine
round trips, some 5\,s over GEO but a third of a second over LEO. Web pages are dominated by such
small transfers, which is why GEO broadband feels sluggish even when its throughput is high,
and why the 25--40\,ms latency of LEO systems was their decisive selling point.

\mdcfig{ch23_tcp_delay}{Left: maximum throughput of a window-limited protocol versus round-trip
time for three window sizes. Right: time to fetch an object with TCP slow start (initial window
10 segments, handshake plus request), assuming unlimited bandwidth: delay, not capacity,
dominates.}

\begin{inpractice}[Performance-enhancing proxies]
GEO broadband operators have long split each TCP connection into three: the user's terminal
and the hub each run a \term{performance-enhancing proxy} (PEP) that terminates TCP locally,
acknowledging data immediately, and carries the traffic across the satellite with a protocol
tuned for a long, clean, high-capacity pipe (large windows, no slow start, selective repeat).
The IETF documented the technique in RFC~3135. PEPs also compress headers and prefetch the
objects a web page will request. Their effectiveness is eroding, however, because the QUIC
transport used by HTTP/3 encrypts its headers, and a proxy can no longer see or spoof its
acknowledgements; operators now rely on tuning at the endpoints and on congestion controllers
that estimate bandwidth and delay directly rather than reacting to loss. Low-orbit systems
largely escape the problem.
\end{inpractice}

\subsection{Doppler shift in low orbit}
\label{sec:ch23:doppler}

\mdcfig{ch23_leo_doppler}{A pass of a 550\,km satellite over a ground station, at 11.7\,GHz, for
three maximum elevations, while the satellite is above 10$^\circ$. Top: elevation and
one-way delay. Bottom: Doppler shift and its rate. Overhead passes last about eight minutes;
the Doppler rate peaks at closest approach. (\texttt{satellite.leo\_pass}.)}

A GEO satellite is almost stationary relative to the ground, so its Doppler shift is tiny,
caused only by its residual station-keeping motion. A LEO satellite moves at 7.6\,km/s, and the
range to it changes at up to about 7\,km/s. Consider the simplest case, a pass directly
overhead, and neglect the Earth's rotation (it changes the answer by a few per cent). If the
satellite moves on a circle of radius $r$ at angular rate $\omega_s=\sqrt{\mu/r^3}$, the central
angle to the station is $\gamma(t)=\omega_st$, measured from the moment of closest approach, and
\begin{equation}
d^2(t)=R_\oplus^2+r^2-2R_\oplus r\cos\gamma(t),\qquad
\dot d=\frac{R_\oplus r\,\omega_s\sin\gamma}{d},\qquad
f_D=-\frac{f_c\,\dot d}{c}.
\label{eq:ch23:doppler}
\end{equation}
The Doppler shift is largest at the horizon, where $\cos\gamma=R_\oplus/r$ and
$d=\sqrt{r^2-R_\oplus^2}$; substituting gives the neat result
$\dot d_{\max}=R_\oplus\omega_s=v\,R_\oplus/r$: the range rate at the horizon is the orbital
speed scaled by $R_\oplus/r$. Differentiating once more at zenith, where $\dot d=0$ and $d=h$,
gives the peak \emph{Doppler rate}:
\begin{equation}
|f_D|_{\max}=\frac{f_c}{c}\,v\,\frac{R_\oplus}{r},\qquad
|\dot f_D|_{\max}=\frac{f_c}{c}\,\frac{v^2R_\oplus}{r\,h}.
\label{eq:ch23:dopplermax}
\end{equation}

\begin{table}[htbp]\centering\small
\caption{Peak Doppler shift and Doppler rate for an overhead pass at 550\,km
($v=7.59$\,km/s), from Equation~\eqref{eq:ch23:dopplermax}.}\label{tab:ch23:doppler}
\begin{tabular}{@{}l r r r@{}}
\toprule
Carrier & $|f_D|_{\max}$ & $|\dot f_D|_{\max}$ & Relative\\
\midrule
700\,MHz (direct-to-cell, low band) & 16.3\,kHz & 225\,Hz/s & 23\,ppm\\
2\,GHz (S-band NTN) & 46.6\,kHz & 642\,Hz/s & 23\,ppm\\
11.7\,GHz (Ku downlink) & 273\,kHz & 3.76\,kHz/s & 23\,ppm\\
20\,GHz (Ka downlink) & 466\,kHz & 6.42\,kHz/s & 23\,ppm\\
\bottomrule
\end{tabular}
\end{table}

For a 550\,km orbit, $v=7.59$\,km/s and $R_\oplus/r=0.921$, so the range rate reaches
$\pm7.0$\,km/s: 23 parts per million of the carrier frequency. Table~\ref{tab:ch23:doppler}
and Figure~\ref{fig:ch23_leo_doppler} give the numbers. At Ku-band the carrier sweeps through
more than half a megahertz in a few minutes, changing fastest---almost 4\,kHz per
second---as the satellite passes overhead. The same range rate also compresses or stretches
the symbol clock by 23\,ppm, so timing recovery must track too, and an NR carrier with
15\,kHz subcarriers at 2\,GHz would be shifted by about three subcarriers. Satellite receivers
handle this with wide acquisition searches and frequency-locked loops aided by ephemeris
(Chapter~\ref{ch:sync}); 3GPP NTN terminals instead \emph{pre-compensate}, using their GNSS
position and the broadcast satellite orbit (Section~\ref{sec:ch23:ntn}).

The same figure shows the other LEO facts of life. A pass above 10$^\circ$ lasts at most
about eight minutes (\texttt{max\_pass\_duration} gives 7.9\,min), and above 25$^\circ$ only
four and a half, so a user terminal must hand over between satellites every few minutes---in
practice more often, because the network reassigns users to balance load and avoid
interference. The one-way delay varies from 1.8 to about 6\,ms during a pass, a changing delay
that timing-advance loops must follow.

\begin{keyidea}[Geometry sets the budget before any electronics]
Altitude and elevation fix the path loss, the delay, the Doppler, the pass duration and the
number of satellites needed for coverage. Going from GEO to a 550\,km LEO cuts the slant range
from about 38\,000\,km to under 1000\,km---more than 30\,dB less spreading loss and a round trip
twenty times shorter---but multiplies the number of satellites by a thousand and turns a fixed
dish into a tracking phased array. Every satellite system design starts with this trade.
\end{keyidea}

\section{Spectrum and regulation}
\label{sec:ch23:spectrum}

\subsection{The frequency bands}

Satellite engineers name bands with the letters inherited from wartime radar, and every band
has its own personality (Table~\ref{tab:ch23:bands}). Lower frequencies propagate through rain
and foliage and suit small, non-directional antennas---hence mobile services in L- and S-band---but
bandwidth there is scarce. Higher frequencies offer gigahertz of spectrum and let a given
antenna form narrower beams, but rain attenuates them heavily. The history of commercial
satellites is a steady climb: C-band in the 1960s, Ku-band for TV from the 1980s, Ka-band for
broadband from the 2000s, and Q/V- and E-band feeder links and optical crosslinks now.

\begin{table}[htbp]\centering\small
\caption{Satellite frequency bands and typical uses. Allocations vary by ITU Region;
values are the common commercial ranges.}\label{tab:ch23:bands}
\begin{tabularx}{\linewidth}{@{}l >{\raggedright\arraybackslash}p{3.7cm} >{\raggedright\arraybackslash}X@{}}
\toprule
Band & Typical satellite allocations & Uses and character\\
\midrule
UHF/VHF & 137--138, 240--400\,MHz & Little-LEO messaging (ORBCOMM), military UHF satcom, data collection\\
L (1--2\,GHz) & 1525--1559 / 1626.5--1660.5\,MHz; 1616--1626.5\,MHz & Mobile satellite service (Inmarsat, Iridium, Thuraya); GNSS at 1.2--1.6\,GHz; ionospheric scintillation\\
S (2--4\,GHz) & 1980--2010 / 2170--2200\,MHz; 2.3\,GHz & MSS and 3GPP NTN band n256; satellite radio (Sirius\,XM); deep-space uplinks (historic)\\
C (4--8\,GHz) & 3.7--4.2 / 5.925--6.425\,GHz & Original ``6/4'' FSS band; TV distribution, trunking; rain-robust but needs large dishes; shared with 5G\\
X (8--12\,GHz) & 7.25--7.75 / 7.9--8.4\,GHz; 8.4--8.45\,GHz & Military satcom; deep-space downlink\\
Ku (12--18\,GHz) & 10.7--12.75 / 13.75--14.5\,GHz & DTH TV, VSAT, maritime/aero, Starlink/OneWeb user links\\
Ka (27--40\,GHz\,$^{*}$) & 17.7--20.2 / 27.5--30\,GHz & HTS broadband, gateways; deep space at 32\,GHz; heavy rain fades\\
Q/V (40--75\,GHz) & 37.5--42.5 / 47.2--51.4\,GHz & HTS feeder links; very heavy rain fades, site diversity\\
E (60--90\,GHz) & 71--76 / 81--86\,GHz & Gateway links for new constellations\\
Optical & 1064, 1550\,nm & Intersatellite laser links, deep-space optical, optical feeder links\\
\bottomrule
\multicolumn{3}{@{}p{\linewidth}@{}}{\footnotesize $^{*}$The satellite industry calls the 17.7--31\,GHz FSS allocations ``Ka-band'' although the IEEE letter band starts at 27\,GHz.}
\end{tabularx}
\end{table}

\begin{inpractice}[The C-band sale]
The 3.7--4.2\,GHz downlink band carried nearly all US cable-TV distribution for four decades,
received by thousands of large dishes at cable head-ends. It is also ideal mid-band spectrum for
5G (Chapter~\ref{ch:cellular}). In 2020 the US FCC ordered satellite operators to vacate
3.7--3.98\,GHz, compensating them for moving their services into 4.0--4.2\,GHz with better
compression, new satellites and filters on every receiving dish; the 2021 auction of the freed
spectrum raised about 81 billion dollars in gross bids. The episode is a lesson in how spectrum
value shifts: a band that was worth little to terrestrial users in 1962 became one of the most
valuable in the world. It also created a new interference problem---high-power 5G next to
sensitive satellite receivers---that filters and guard bands must solve.
\end{inpractice}

\subsection{The ITU and the orbital ring}

Radio waves ignore borders, and a geostationary satellite is visible from dozens of countries,
so satellite spectrum is governed internationally by the \term{Radio Regulations} of the
International Telecommunication Union, revised at a World Radiocommunication Conference (WRC)
every three to four years. The Regulations allocate bands to \emph{services}---fixed-satellite
(FSS), broadcasting-satellite (BSS), mobile-satellite (MSS) and others---in each of three
world regions. An operator's national administration files the planned network with the ITU,
coordinates it with every existing or earlier-filed network that it might interfere with, and
finally notifies it for recording in the Master International Frequency Register, which gives
international recognition and protection. For most FSS bands the principle is ``first come,
first served'', which encourages speculative ``paper satellites''; for some BSS and FSS bands,
plans drawn up at conferences since 1977 guarantee each country orbital positions and
spectrum whether or not it can yet use them. For non-geostationary constellations,
WRC-19 introduced milestones: an operator must deploy set percentages of its filed
constellation within set periods or see its filing reduced.

In the geostationary ring, neighbouring satellites using the same band are typically spaced
2--3$^\circ$ apart. The limiting factor is the \emph{earth station} antenna: a small dish has a
wide beam and sidelobes that illuminate neighbouring satellites on the uplink and receive
their signals on the downlink. ITU-R Recommendation S.580 sets a design objective that
earth-station sidelobe gain stay below $29-25\log_{10}\theta$\,dBi for off-axis angles $\theta$
from 1$^\circ$ to 20$^\circ$ (Figure~\ref{fig:ch23_antennas}), and operators and regulators
limit the off-axis EIRP density of small uplink terminals. The adjacent-satellite interference
that remains is one of the terms in the total $C/(N+I)$ of the link
(Section~\ref{sec:ch23:combining}).

\paragraph{GSO and NGSO sharing.} When thousands of non-geostationary (NGSO) satellites share
Ku- and Ka-band with geostationary networks, a new geometry arises: whenever a LEO satellite
passes between a GEO earth station and its satellite, it is in the main beam of the earth
station. Article~22 of the Radio Regulations therefore sets limits on the \term{equivalent
power flux density} (EPFD)---the aggregate power flux density from all satellites of an NGSO
system at any point on the Earth, weighted by the receiving antenna's gain towards each of them
relative to its peak---which NGSO systems must not exceed at GEO earth stations. In practice LEO
systems implement \emph{GSO arc avoidance}: a satellite stops transmitting to (or a user terminal
stops pointing at) any direction within some angle of the geostationary arc as seen from the
ground. This is one reason LEO terminals at low latitudes look away from the equator, and one
reason LEO capacity in equatorial regions is harder to deliver.

\section{The satellite link budget}
\label{sec:ch23:linkbudget}

Chapter~\ref{ch:noise} introduced the link budget. Here we develop it with the vocabulary and
precision of the satellite industry, where a single decibel may be worth millions of dollars of
satellite power or antenna size.

\subsection{The fundamental equation}

A transmitter of power $P_t$ feeding an antenna of gain $G_t$ produces an \term{effective
isotropic radiated power} $\text{EIRP}=P_tG_t$ (after any feed losses). At range $d$ the power
flux density is $\text{EIRP}/4\pi d^2$ W/m$^2$, and a receiving antenna of effective area
$A_e=G_r\lambda^2/4\pi$ collects $C=\text{EIRP}\,G_r/L_{\text{fs}}$, with
$L_{\text{fs}}=(4\pi d/\lambda)^2$. Dividing by the noise density $N_0=kT_{\text{sys}}$ gives the
central equation of satellite engineering:
\begin{equation}
\boxed{\;\frac{C}{N_0}=\text{EIRP}-L_{\text{fs}}-L_a-L_{\text{misc}}+\frac{G}{T}-k
\quad\text{[dB-Hz]}\;}
\label{eq:ch23:cn0}
\end{equation}
where $L_a$ is the atmospheric loss, $L_{\text{misc}}$ collects pointing, polarisation and
feed losses, $G/T=G_r/T_{\text{sys}}$ is the receiver's \term{figure of merit} in dB/K, and
$k=-228.6$\,dBW/K/Hz, so subtracting $k$ adds 228.6\,dB. The equation separates cleanly into a
transmitter term (EIRP), a path term, and a receiver term ($G/T$), which is why satellites are
specified by their EIRP and $G/T$ contours over the coverage area and terminals by their
EIRP and $G/T$. From $C/N_0$ one obtains
$E_b/N_0=C/N_0-10\log_{10}R_b$ and, for a carrier of symbol rate $R_s$,
$E_s/N_0=C/N_0-10\log_{10}R_s$, which (with an ideal matched filter) is also the carrier-to-noise
ratio $C/N$ in the noise bandwidth $R_s$. The functions \texttt{cn0\_dbhz}, \texttt{fspl\_db},
\texttt{dish\_gain\_dbi} and the \texttt{LinkBudget} class in \texttt{commlib/satellite.py}
implement all of this.

\subsection{Receiver noise: $G/T$}

The system noise temperature of an earth station, referred to the LNA input, is
\begin{equation}
T_{\text{sys}}=\frac{T_{\text{ant}}}{L_f}+T_0\left(1-\frac{1}{L_f}\right)+T_{\text{LNA}}+\frac{T_{\text{rest}}}{G_{\text{LNA}}} ,
\label{eq:ch23:tsys}
\end{equation}
where $L_f$ is the loss of any feed components ahead of the LNA at physical temperature $T_0$
(Chapter~\ref{ch:noise}). The antenna temperature $T_{\text{ant}}$ is the sky brightness
temperature weighted by the antenna pattern: at Ku-band and 30$^\circ$ elevation in clear
air the sky contributes perhaps 10--20\,K and ground pickup through spillover and sidelobes
another 10--30\,K. A modern Ku LNB has a noise figure of 0.5--1\,dB (35--75\,K), so a DTH
terminal's $T_{\text{sys}}$ is around 80--120\,K. Every 0.1\,dB of loss in front of the LNB adds
about 7\,K, which is why the LNB is bolted directly to the feed horn.

A satellite's receiving antenna looks down at the 290\,K Earth, so its $T_{\text{ant}}$ is
around 290\,K and its system temperature is typically 500--1000\,K; satellite $G/T$ values range
from about $-10$\,dB/K for a global beam to $+20$\,dB/K or more for a small HTS spot beam.

\subsection{Atmospheric impairments}

Between the satellite and the ground, the signal crosses the ionosphere and the troposphere.
Below about 3\,GHz the \emph{ionosphere} matters: its free electrons rotate the plane of linear
polarisation (Faraday rotation) and, near the geomagnetic equator after sunset, cause rapid
amplitude and phase \emph{scintillation}. That is why L- and S-band mobile systems use circular
polarisation. Above about 10\,GHz the \emph{troposphere} dominates:
\begin{itemize}
\item \textbf{Gaseous absorption} (ITU-R P.676) by oxygen and water vapour: a few tenths of a
decibel at Ku- and Ka-band on typical slant paths, with a water-vapour line at 22.2\,GHz and a
wall of oxygen absorption around 60\,GHz (used deliberately for crosslinks that cannot be
heard from the ground).
\item \textbf{Clouds and fog} (ITU-R P.840): a fraction of a decibel at Ku, up to a few decibels
at Q/V-band.
\item \textbf{Tropospheric scintillation}: fast fluctuations of a fraction of a decibel, larger at
low elevation.
\item \textbf{Rain} (ITU-R P.618, P.837, P.838, P.839): by far the largest effect above 10\,GHz.
\end{itemize}

\paragraph{Rain.} Raindrops absorb and scatter microwaves when their size becomes a
significant fraction of the wavelength. Empirically, the specific attenuation depends on the
rain rate $R$ (mm/h) as a power law,
\begin{equation}
\gamma_R=kR^{\alpha}\quad\text{dB/km},
\label{eq:ch23:krate}
\end{equation}
with frequency- and polarisation-dependent coefficients tabulated in ITU-R P.838
(Figure~\ref{fig:ch23_rain}, left). At 12\,GHz, a 50\,mm/h downpour attenuates about
2.2\,dB/km; at 30\,GHz, about 9\,dB/km; at 4\,GHz, a few hundredths. Rain is not uniform,
however: heavy rain falls in cells a few kilometres across, and only the part of the path
below the \emph{rain height} (roughly the 0$^\circ$C isotherm, 3--5\,km at mid-latitudes) is
wet. The ITU-R P.618 prediction method, used worldwide for satellite system design,
turns a single climatic statistic into an attenuation distribution:
\begin{enumerate}
\item Find the rain height $h_R$ (P.839) and the point rain rate $R_{0.01}$ exceeded 0.01\% of an
average year at the site (P.837 maps: roughly 20--40\,mm/h in temperate climates and 60--150\,mm/h
in the tropics).
\item Compute the slant path below the rain height, $L_s=(h_R-h_s)/\sin\theta$, and its horizontal
projection $L_G=L_s\cos\theta$.
\item Compute $\gamma_R$ from~\eqref{eq:ch23:krate} with $R=R_{0.01}$.
\item Apply a \emph{horizontal reduction factor} and a \emph{vertical adjustment factor}, empirical
functions of $L_G\gamma_R/f$, the elevation and latitude that account for the finite size of rain
cells, to get an effective path length $L_E$. Then $A_{0.01}=\gamma_RL_E$.
\item Scale to other percentages $p$ with
\[
A_p=A_{0.01}\left(\frac{p}{0.01}\right)^{-(0.655+0.033\ln p-0.045\ln A_{0.01}-\beta(1-p)\sin\theta)},
\]
where $\beta$ is a latitude-dependent correction.
\end{enumerate}
The function \texttt{rain\_atten\_p618} follows these steps. Figure~\ref{fig:ch23_rain} (right)
shows the result for a temperate site with $R_{0.01}=42$\,mm/h and a 35$^\circ$ path. At
12\,GHz, a margin of 3\,dB buys 99.9\% availability and 8\,dB buys 99.99\%; at 30\,GHz the same
availabilities need about 17 and 44\,dB. No practical fixed margin covers the tail at Ka-band and
above, which is why those systems use adaptive coding and modulation, uplink power control and,
for gateways, \emph{site diversity}---two gateways tens of kilometres apart are rarely under the
same rain cell.

\mdcfig{ch23_rain}{Left: specific attenuation of rain versus frequency for five rain rates
(coarse ITU-R P.838-3 coefficients, circular polarisation). Right: slant-path rain attenuation
exceeded for a given percentage of an average year at four frequencies (35$^\circ$ elevation,
$R_{0.01}=42$\,mm/h, latitude 45$^\circ$, following the ITU-R P.618 procedure). The top axis
converts the percentage to availability.}

\begin{keyidea}[Rain hurts a downlink twice]
An absorbing medium also emits. A rain layer with attenuation $A$ at physical temperature
$T_m\approx275$\,K adds a sky noise temperature
\begin{equation}
\Delta T=T_m\left(1-10^{-A/10}\right)
\label{eq:ch23:rainnoise}
\end{equation}
to the antenna. A 3\,dB fade adds about 140\,K---more than the entire clear-sky system
temperature of a good Ku-band terminal. On a low-noise downlink, the rise in noise can cost as
much as the attenuation itself: the total degradation in $C/N$ is $A+10\log_{10}(T_{\text{sys}}'/T_{\text{sys}})$.
The uplink escapes the second penalty, because the satellite already looks at a 290\,K Earth. The
quieter your receiver, the more rain hurts in relative terms.
\end{keyidea}

Two other rain effects matter. Non-spherical raindrops and ice crystals \emph{depolarise} the
signal, reducing the cross-polar discrimination on which dual-polarisation frequency reuse
depends. And rain statistics are for an \emph{average year}; the worst month is considerably
worse (ITU-R P.841 converts between the two), which matters for service-level agreements.

\subsection{Uplink, downlink and interference}
\label{sec:ch23:combining}

In a \term{bent-pipe} (transparent) satellite, the transponder amplifies whatever it receives,
noise included. Let the uplink deliver carrier $C_u$ and noise $N_u$ to the satellite, and the
downlink path have overall gain $g$ from the satellite input to the ground receiver, where the
receiver adds its own noise $N_d$. At the ground the carrier is $gC_u$ and the noise is
$gN_u+N_d$, so
\begin{equation}
\left(\frac{C}{N}\right)_{\text{tot}}^{-1}=\frac{gN_u+N_d}{gC_u}=
\left(\frac{C}{N}\right)_{u}^{-1}+\left(\frac{C}{N}\right)_{d}^{-1},
\end{equation}
with $(C/N)_d=gC_u/N_d$ the downlink ratio. Interference that is noise-like adds in the same
way, so the general rule is
\begin{equation}
\left(\frac{C}{N+I}\right)_{\text{tot}}^{-1}=\sum_i\left(\frac{C}{X_i}\right)^{-1},
\qquad X_i\in\{N_u,\,N_d,\,I_{\text{ASI}},\,I_{\text{XPOL}},\,I_{\text{IM}},\,I_{\text{CCI}},\dots\}
\label{eq:ch23:combine}
\end{equation}
over uplink and downlink noise, adjacent-satellite interference, cross-polar interference,
intermodulation in the transponder and co-channel interference from other beams. The total is
always worse than the worst contribution, and when two terms are equal they cost 3\,dB. A
well-designed link is \emph{balanced}: no term is so good that its cost (in satellite power or
antenna size) is wasted, and none so bad that it dominates. \texttt{combine\_cn\_db} implements
the rule.

A \term{regenerative} payload demodulates and decodes the uplink and retransmits clean
symbols. Then noise does not accumulate; instead the error probabilities approximately add,
$P_e\approx P_{e,u}+P_{e,d}$. Because error probability falls so steeply with SNR, the total is
dominated by the weaker link, and a regenerative satellite typically needs a few decibels less
than a transparent one for balanced links---one of several reasons for on-board processing
(Section~\ref{sec:ch23:payload}).

\subsection{Worked link budgets}

\begin{worked}[A GEO Ku-band DTH downlink]
A broadcast satellite at 19.2$^\circ$E delivers 51\,dBW EIRP towards a viewer near Paris
(near the edge of its main coverage) at 11.7\,GHz. The geometry of the previous worked
example gives an elevation of about 32$^\circ$ and a range of 38\,400\,km, so
$L_{\text{fs}}=20\log_{10}(4\pi\cdot3.84\times10^7\cdot1.17\times10^{10}/3\times10^8)=205.5$\,dB.
Allow 0.8\,dB for clear-sky gases, pointing and polarisation. The 60\,cm dish with 65\%
efficiency has $G=10\log_{10}[0.65(\pi\cdot0.6/0.0256)^2]=35.5$\,dBi. With an antenna temperature
of 35\,K and an LNB of 0.7\,dB noise figure (51\,K), $T_{\text{sys}}=86$\,K $=19.3$\,dBK, so
$G/T=16.1$\,dB/K. Then
\[
C/N_0=51-205.5-0.8+16.1+228.6=89.4\ \text{dB-Hz}.
\]
For a 27.5\,Mbaud carrier, $C/N=89.4-74.4=15.0$\,dB. Now combine it with the other contributions
of~\eqref{eq:ch23:combine}: uplink $C/N$ 25\,dB, adjacent satellites 21\,dB, cross-polarisation
22\,dB, transponder intermodulation 20\,dB. The total is 12.3\,dB (Figure~\ref{fig:ch23_dth_budget})---about
2.7\,dB below the downlink alone. Interference is not a detail in satellite links; it is a third
of the budget.

With a DVB-S2 carrier using 8PSK rate 2/3, which needs 6.6\,dB (Table~\ref{tab:ch23:dvbs2}), the
clear-sky margin is 5.7\,dB. How often is it used up? Table~\ref{tab:ch23:dthrain} applies the
P.618 method with $R_{0.01}\approx30$\,mm/h, typical of western Europe, and includes the rise in
sky noise. The link holds for about 99.95\% of the year, failing roughly four hours a year in
the heaviest showers. Notice how, in the 0.1\% row, 2.1\,dB of attenuation costs 5.6\,dB of
$C/N$---the noise penalty of~\eqref{eq:ch23:rainnoise} is larger than the fade itself.
\end{worked}

\mdcfig{ch23_dth_budget}{The GEO Ku-band DTH worked example. Left: the downlink $C/N_0$ built up
term by term. Right: combining the downlink with uplink noise, adjacent-satellite, cross-polar
and intermodulation interference per Equation~\eqref{eq:ch23:combine}.}

\begin{table}[htbp]\centering\small
\caption{Rain on the Ku-band DTH link ($R_{0.01}=30$\,mm/h, 32$^\circ$ elevation, 11.7\,GHz).
``Total'' combines the faded downlink with the fixed uplink and interference terms.}\label{tab:ch23:dthrain}
\begin{tabular}{@{}r r r r r l@{}}
\toprule
Time exceeded & Availability & Rain loss & $T_{\text{sys}}$ & Total $C/(N+I)$ & 8PSK 2/3 (6.6\,dB)?\\
\midrule
clear sky & --- & 0 & 86\,K & 12.3\,dB & yes, 5.7\,dB margin\\
1\% & 99\% & 0.5\,dB & 114\,K & 11.3\,dB & yes\\
0.5\% & 99.5\% & 0.8\,dB & 131\,K & 10.7\,dB & yes\\
0.1\% & 99.9\% & 2.1\,dB & 192\,K & 8.5\,dB & yes\\
0.05\% & 99.95\% & 3.1\,dB & 226\,K & 7.1\,dB & just\\
0.01\% & 99.99\% & 6.6\,dB & 300\,K & 2.8\,dB & no\\
\bottomrule
\end{tabular}
\end{table}

\begin{worked}[A Ka-band HTS spot beam with adaptive coding and modulation]
A high-throughput GEO satellite serves a user at 40$^\circ$N, 30$^\circ$ of longitude away,
from a spot beam with 58\,dBW EIRP at 19.7\,GHz. The elevation is 34.3$^\circ$ and the range
38\,240\,km, so $L_{\text{fs}}=210.0$\,dB. A 75\,cm dish gives 41.9\,dBi; with 60\,K antenna
temperature, a 0.3\,dB feed loss and a 1.5\,dB-noise-figure LNB, $T_{\text{sys}}=195$\,K and
$G/T=19.0$\,dB/K. With 0.8\,dB of miscellaneous losses,
$C/N_0=58-210.0-0.8+19.0+228.6=94.8$\,dB-Hz, and for a 200\,Mbaud carrier $C/N=11.8$\,dB. Co-channel
interference from other beams using the same colour (Section~\ref{sec:ch23:antennas}) is around
20\,dB, giving 11.2\,dB in clear sky: enough for 16APSK rate 3/4 at 2.97\,b/s/Hz with 0.5\,dB margin.

But this site is wetter ($R_{0.01}=42$\,mm/h) and Ka-band fades are deep. Running the P.618 model
and the ACM selector \texttt{acm\_select} of \texttt{commlib.satellite} gives
Table~\ref{tab:ch23:kaacm}. A \emph{constant} coding and modulation (CCM) design sized for 99.9\%
availability would have to use QPSK 2/5 (0.79\,b/s/Hz) all the time. With \emph{adaptive} coding
and modulation the terminal reports its $E_s/N_0$ and the hub picks the best MODCOD frame by frame:
the long-term average efficiency is about 2.93\,b/s/Hz, 3.7 times the CCM design, while
availability is unchanged. This single idea is why every broadband satellite system built since
2005 uses ACM.
\end{worked}

\begin{table}[htbp]\centering\small
\caption{Ka-band ACM example: 19.7\,GHz, 34$^\circ$ elevation, $R_{0.01}=42$\,mm/h, 0.5\,dB
margin, DVB-S2 MODCODs.}\label{tab:ch23:kaacm}
\begin{tabular}{@{}r r r r l r@{}}
\toprule
Time exceeded & Rain loss & $T_{\text{sys}}$ & $C/(N+I)$ & MODCOD & b/s/Hz\\
\midrule
clear sky & 0 & 195\,K & 11.2\,dB & 16APSK 3/4 & 2.97\\
5\% & 0.7\,dB & 231\,K & 10.0\,dB & 16APSK 2/3 & 2.64\\
1\% & 2.1\,dB & 293\,K & 7.7\,dB & 8PSK 2/3 & 1.98\\
0.5\% & 3.3\,dB & 331\,K & 6.1\,dB & 8PSK 3/5 & 1.78\\
0.1\% & 8.2\,dB & 412\,K & 0.4\,dB & QPSK 2/5 & 0.79\\
0.05\% & 11.5\,dB & 433\,K & $-3.2$\,dB & outage (DVB-S2) & 0\\
\bottomrule
\end{tabular}
\end{table}

\begin{worked}[A LEO user terminal]
Operators of LEO constellations do not publish complete link budgets, so take this one as an
illustration with plausible, rounded assumptions. A satellite at 550\,km serves a user at
40$^\circ$ elevation (slant range 812\,km) at 11.7\,GHz with a beam EIRP of 36\,dBW spread over a
240\,MHz channel. The free-space loss is only 172.0\,dB---33\,dB less than from GEO. The user
terminal is a flat phased array of about 0.23\,m$^2$; with 70\% aperture efficiency its broadside
gain is $10\log_{10}(4\pi\cdot0.23\cdot0.7/\lambda^2)=34.9$\,dBi, and with an assumed
$T_{\text{sys}}=150$\,K (no cryogenics, a few tenths of a decibel of loss in the beamformer
before each LNA) its $G/T$ is 13.1\,dB/K. Scanning 30$^\circ$ off broadside costs about 0.8\,dB
(Section~\ref{sec:ch23:antennas}). Allowing 0.5\,dB of atmospheric loss,
\[
C/N_0=36-172.0-0.5+13.1-0.8+228.6=104.4\ \text{dB-Hz},
\]
so $C/N=104.4-83.8=20.6$\,dB in 240\,MHz. The Shannon limit at that SNR is about 6.9\,b/s/Hz,
some 1.6\,Gb/s for the channel---shared, of course, among all the users in the beam. The lesson
is the one in the geometry key idea: compared with the GEO DTH example, the user antenna has
about the same gain and a worse noise temperature, and the satellite radiates 15\,dB less
EIRP---yet the 33\,dB saved in path loss still leave an SNR high enough for gigabit channels.
The price is paid elsewhere: in the number of satellites, in electronic beam steering, and in
handovers every few minutes.
\end{worked}

\section{Antennas and beams}
\label{sec:ch23:antennas}

\subsection{Reflector antennas}

The parabolic reflector is the workhorse of satellite communication: it converts the
spherical wave from a small feed at its focus into a plane wave across its aperture. Its gain is
\begin{equation}
G=\eta\left(\frac{\pi D}{\lambda}\right)^2,\qquad
\theta_{3\text{dB}}\approx 70\,\frac{\lambda}{D}\ \text{degrees},
\label{eq:ch23:dishgain}
\end{equation}
where $\eta$, the aperture efficiency, is typically 0.55--0.75. It is the product of several
factors: the \emph{illumination efficiency}, lost because the feed tapers the illumination
towards the rim to keep sidelobes low; the \emph{spillover} of feed energy past the rim (which
also picks up warm ground noise); \emph{blockage} by the feed and struts; and the
\emph{surface accuracy}, captured by Ruze's formula $\eta_s=\exp[-(4\pi\varepsilon/\lambda)^2]$
for an rms surface error $\varepsilon$. A 0.5\,mm rms error costs 0.3\,dB at 12\,GHz but 3\,dB
at 40\,GHz---which is why Ka- and Q-band antennas are expensive and why the DSN's 70\,m dishes
needed painstaking refurbishment to work well at 32\,GHz.

The beamwidth sets the \emph{pointing} requirement. A 60\,cm dish at 11.7\,GHz has
$\theta_{3\text{dB}}\approx3^\circ$; a pointing error $\Delta\theta$ costs approximately
$12(\Delta\theta/\theta_{3\text{dB}})^2$\,dB, so 0.3$^\circ$ costs 0.12\,dB. A 9\,m gateway antenna at
30\,GHz has a beamwidth of about 0.08$^\circ$ and must actively track even a well-kept GEO
satellite, using \emph{step-tracking} (small deliberate mispointings to find the peak) or
\emph{monopulse} feeds that measure the pointing error directly.

Most consumer dishes are \term{offset-fed}: the reflector is a section of a paraboloid cut away
from its axis, so the feed and its support sit outside the beam and cause no blockage---which is
why a DTH dish looks oval and seems to point at the ground. Large earth stations use
\emph{Cassegrain} or \emph{Gregorian} dual-reflector systems, whose subreflector folds the optics
so that heavy electronics sit behind the main reflector; the DSN antennas go further and guide
the beam through a series of mirrors into a room below, where cryogenic receivers stay still while
the dish moves. The feed itself is usually a corrugated horn, whose grooves give a symmetric beam
with very low cross-polarisation, followed by an \emph{orthomode transducer} (OMT) that separates
two orthogonal polarisations---and so doubles the spectrum a link can use.

\mdcfig{ch23_antennas}{Left: patterns of 0.6, 1.2 and 2.4\,m reflectors at 14.25\,GHz (tapered
circular apertures) against the ITU-R S.580 sidelobe objective $29-25\log_{10}\theta$. Small
dishes put significant gain towards a neighbouring satellite 2$^\circ$ away. Right: a 32-element
linear phased array with half-wavelength spacing, scanned to 0$^\circ$, 30$^\circ$ and 55$^\circ$:
the beam broadens and its gain falls as it scans.}

Figure~\ref{fig:ch23_antennas} (left) shows why small dishes are a regulatory problem. A 60\,cm
dish's main lobe is still several decibels wide at 2$^\circ$, where the next satellite sits;
its uplink therefore radiates into its neighbour, and its downlink hears the neighbour. The
remedies are to limit the off-axis EIRP density of small terminals (spreading their signals, if
necessary), to coordinate frequency plans with neighbours, and to accept the resulting
adjacent-satellite interference as a term in the budget.

\subsection{Spot beams and frequency reuse}

A satellite antenna of diameter $D$ forms a beam whose footprint on the ground has a diameter of
about $d\,\theta_{3\text{dB}}$. At Ka-band a 2.6\,m reflector gives a beam about 0.4$^\circ$ wide:
from GEO, a spot some 250\,km across. Covering the continental United States then takes on the
order of a hundred beams---and each beam can \emph{reuse the spectrum}. That idea, borrowed from
the cellular concept (Chapter~\ref{ch:multipleaccess}), created the high-throughput satellite.

The usual pattern is \term{four-colour reuse} (Figure~\ref{fig:ch23_beam_reuse}): the available
band is split into two halves and each is used on two orthogonal polarisations, giving four
``colours'' arranged so that adjacent beams never share one. The co-channel beams are then about
two beam spacings apart, where the antenna pattern is 20--30\,dB down; the aggregate co-channel
interference from all same-colour beams typically leaves a $C/I$ of 15--20\,dB at the beam
edges---one of the dominant terms in an HTS budget.

\mdcfig{ch23_beam_reuse}{Left: four-colour frequency reuse over a hexagonal lattice of spot beams
(two sub-bands times two polarisations). Right: cuts through neighbouring beams along a row; the
nearest beams of the same colour (navy) are two beam spacings away, well down the pattern of the
wanted beam.}

The capacity arithmetic is simple and striking. Suppose the Ka-band user downlink has 2.5\,GHz of
spectrum on each of two polarisations, 5\,GHz in all. With four colours, each beam gets 1.25\,GHz.
A satellite with 100 beams therefore radiates 125\,GHz of total bandwidth; at an average of
2.5\,b/s/Hz under ACM, that is over 300\,Gb/s---from one satellite, with the spectrum of a single
wide-beam system. Doubling the number of beams (a bigger antenna, smaller beams) doubles the
capacity again, until power, mass, heat and above all the \emph{feeder links} run out: every
bit sent to users must first come up from a gateway. A gateway using the same 5\,GHz on its
uplink can feed only a few user beams' worth of traffic, so a large HTS needs dozens of
gateways, spread far enough apart to reuse their own spectrum and placed where rain is least
likely. Moving feeder links to Q/V-band (with several gigahertz more spectrum per gateway) or to
optical links is an active area of development.

Uniform four-colour reuse wastes capacity where there are few users and starves the busy
beams. Two refinements address this. \term{Beam hopping} illuminates each beam only during the
time slots it needs, so a single amplifier and band can serve several beams in turn according
to demand; DVB-S2X's superframe format was designed partly to support it. \emph{Multibeam
precoding}, the satellite analogue of multi-user MIMO (Chapter~\ref{ch:mimo}), goes further:
with channel knowledge from the terminals, the gateway precodes the signals of neighbouring beams
so that their interference cancels at each user, allowing full-frequency reuse in every beam.

\subsection{Phased arrays for moving beams}

A LEO satellite moves across the sky in minutes, so its users need antennas that track---and
that must switch to a new satellite at the other side of the sky in milliseconds. A
mechanically steered dish can do the first, but not the second at consumer cost. The answer is
the \term{phased array}: a flat panel of many small antenna elements, each with its own phase
shifter (or time delay) and amplifier, whose beam is steered electronically by setting the
phases. For an array with element spacing $s$ steered to angle $\theta_0$, element $n$ is given
the phase $-2\pi ns\sin\theta_0/\lambda$; the array factor of an $N$-element line is then
\begin{equation}
|AF(\theta)|=\left|\frac{\sin[N\pi s(\sin\theta-\sin\theta_0)/\lambda]}{N\sin[\pi s(\sin\theta-\sin\theta_0)/\lambda]}\right| .
\end{equation}
Three facts follow (Figure~\ref{fig:ch23_antennas}, right). The spacing must be at most about
$\lambda/2$ to avoid \emph{grating lobes}---copies of the main beam---at wide scan angles, so a
panel at 12\,GHz, where $\lambda/2=1.3$\,cm, has roughly 600 elements per 0.1\,m$^2$. The beam
widens by $1/\cos\theta_0$ as it scans. And the gain falls, because the projected aperture
shrinks as $\cos\theta_0$ and the element pattern also rolls off: \emph{scan loss} is roughly
$10x\log_{10}(1/\cos\theta_0)$\,dB with $x\approx1.2$--1.5, about 1\,dB at 30$^\circ$ and 3--4\,dB at
55$^\circ$ (\texttt{array\_scan\_loss\_db}). LEO user terminals are therefore often mounted
tilted, and constellations prefer to serve users at high elevation.

The engineering revolution behind LEO broadband is that such arrays---once the preserve of
military radars costing millions---are now built from silicon beamformer chips, each handling
a few elements with integrated phase shifters, amplifiers and control, on ordinary printed
circuit boards, for a few hundred dollars. The same technology steers the beams of the
satellites themselves, and of aircraft terminals that must look past the wing at low
elevation. Digital beamforming, in which every element (or subarray) has its own converter and
the beams are formed in DSP (Chapter~\ref{ch:dsp}), lets one array form many independent beams
at once and is used on the newest satellites.

\section{Inside the satellite: the payload}
\label{sec:ch23:payload}

\subsection{The bent-pipe transponder}

\begin{figure}[htbp]\centering
\begin{tikzpicture}[x=1cm,y=1cm,
  blk/.style={draw=navy,fill=navy!6,thick,minimum height=0.8cm,minimum width=1.1cm,align=center,font=\sffamily\scriptsize},
  amp/.style={draw=navy,fill=accent!10,thick,isosceles triangle,isosceles triangle apex angle=60,minimum height=0.62cm,minimum width=0.62cm,inner sep=0pt},
  arr/.style={-{Stealth[length=1.8mm]},thick,navy}]
% receive side
\draw[navy,thick] (-0.35,0.45) -- (0,0) -- (-0.35,-0.45); \draw[navy,thick] (-0.35,0.45) to[bend right=40] (-0.35,-0.45);
\node[font=\sffamily\scriptsize,navy,align=center] at (-0.3,0.85) {Rx\\antenna};
\node[font=\sffamily\scriptsize,accent] at (-0.2,-0.75) {6\,GHz};
\node[blk] (bpf) at (1.0,0) {input\\filter};
\node[amp] (lna) at (2.25,0) {};
\node[font=\sffamily\scriptsize,above=0.12cm of lna,align=center] {LNA};
\node[draw=navy,thick,circle,minimum size=0.55cm,font=\small] (mix) at (3.3,0) {$\times$};
\node[blk,fill=white,minimum width=1.0cm] (lo) at (3.3,-1.35) {LO\\2225\,MHz};
\node[blk,minimum height=2.9cm,minimum width=1.0cm] (imux) at (4.6,0) {I\\M\\U\\X};
\draw[arr] (0,0) -- (bpf); \draw[arr] (bpf) -- (lna); \draw[arr] (lna) -- (mix); \draw[arr] (lo) -- (mix); \draw[arr] (mix) -- (imux);
% channel chains
\foreach \y/\n in {1/1,0/2,-1/3}{
  \node[blk,minimum height=0.62cm,minimum width=1.25cm] (camp\n) at (6.3,\y) {CAMP + lin.};
  \node[amp] (twt\n) at (7.75,\y) {};
  \draw[arr] (imux.east |- camp\n) -- (camp\n); \draw[arr] (camp\n) -- (twt\n);
}
\node[font=\sffamily\scriptsize] at (7.7,1.5) {TWTA};
\node[blk,minimum height=2.9cm,minimum width=1.0cm] (omux) at (9.1,0) {O\\M\\U\\X};
\foreach \n in {1,2,3}{\draw[arr] (twt\n.east) -- (omux.west |- twt\n);}
\node[font=\sffamily\scriptsize,gray] at (6.55,-1.6) {one chain per 36\,MHz channel};
\draw[arr] (omux) -- (10.25,0);
\draw[navy,thick] (10.6,0.45) -- (10.25,0) -- (10.6,-0.45); \draw[navy,thick] (10.6,0.45) to[bend left=40] (10.6,-0.45);
\node[font=\sffamily\scriptsize,navy,align=center] at (10.55,0.85) {Tx\\antenna};
\node[font=\sffamily\scriptsize,accent] at (10.45,-0.75) {4\,GHz};
\end{tikzpicture}
\caption{A classic C-band bent-pipe transponder. The wideband receiver amplifies the whole
500\,MHz uplink band and translates it down by 2225\,MHz; the input multiplexer (IMUX) splits it
into channels, each with its own channel amplifier (CAMP), lineariser and high-power amplifier;
the output multiplexer (OMUX) recombines them for the transmit antenna. Redundancy switching
rings (not shown) let spare amplifiers replace failed ones.}
\label{fig:ch23_transponder}
\end{figure}

The simplest payload, and still the most common, is the \term{transparent} or
\term{bent-pipe} transponder of Figure~\ref{fig:ch23_transponder}. It receives the uplink,
amplifies it with a low-noise amplifier, shifts it to the downlink band with a single mixer
(for the 6/4\,GHz C-band pair, the local oscillator is 2225\,MHz), separates it into
\emph{channels} with an input multiplexer, amplifies each channel with its own high-power
amplifier and recombines them in an output multiplexer. The word \emph{transponder} originally
meant one of these channel chains. A classic C-band satellite divides its 500\,MHz into twelve
channels of 36\,MHz with 4\,MHz guard bands and uses both polarisations to get 24 transponders;
Ku-band transponders of 27, 33, 36, 54 and 72\,MHz are common. The payload does not care what
modulation passes through it---analog FM TV in 1970, DVB-S2 today---which is why bent-pipe
satellites routinely outlive the standards they were designed for.

Why split into channels at all? Because the high-power amplifiers are the payload's most
nonlinear, most power-hungry and least reliable parts. Giving each channel its own amplifier
limits intermodulation to within the channel, lets each amplifier run near its optimum drive,
and lets redundancy rings swap in spares. The multiplexers are high-$Q$ waveguide or dielectric
resonator filters whose sharp skirts inevitably bring amplitude ripple and group-delay
variation near the band edges---linear distortions that the ground demodulator must equalise
(Chapter~\ref{ch:equalization}).

\subsection{High-power amplifiers and their nonlinearity}

Two kinds of amplifier fly. The \term{travelling-wave tube amplifier} (TWTA) passes an electron
beam through a helix along which the RF wave travels at nearly the same speed, transferring
energy from beam to wave. TWTAs deliver 100--300\,W at Ku- and Ka-band with DC-to-RF efficiencies
around 60--70\% and long space heritage. \textbf{Solid-state power amplifiers}\index{solid-state power amplifier} (SSPAs), now mostly
gallium nitride, dominate at L-, S- and C-band, in phased arrays (where each element needs only
watts), and increasingly at higher frequencies. Both saturate: as the drive increases, the output
power approaches a maximum. A TWTA also shifts the phase of the output as the amplitude grows,
because the electrons slow down as they give up more energy. These two effects are described by
the \term{AM/AM} and \term{AM/PM} characteristics. Adel Saleh's classic 1981 model captures a
TWTA with four parameters: for input amplitude $r$,
\begin{equation}
A(r)=\frac{\alpha_ar}{1+\beta_ar^2},\qquad \Phi(r)=\frac{\alpha_\phi r^2}{1+\beta_\phi r^2},
\label{eq:ch23:saleh}
\end{equation}
with the fitted values $\alpha_a=2.1587$, $\beta_a=1.1517$, $\alpha_\phi=4.0033$\,rad,
$\beta_\phi=9.1040$. Saturation occurs at $r=1/\sqrt{\beta_a}$. Figure~\ref{fig:ch23_twta}
shows both curves and, for comparison, the Rapp model of an SSPA, which has negligible AM/PM but
a sharper knee.

\mdcfig{ch23_twta}{Left: AM/AM characteristics of a TWTA (Saleh model) and an SSPA (Rapp model
with $p=3$), normalised to saturation. Right: the TWTA's AM/PM conversion, about 22$^\circ$ at
saturation. (\texttt{satellite.saleh\_twta}, \texttt{commlib.rapp\_pa}.)}

The operating point is described by the \term{input back-off} (IBO), the ratio of the saturating
input power to the actual mean input power, and the \term{output back-off} (OBO), the
corresponding ratio at the output. Operating at saturation extracts all the power the satellite
has paid for, but distorts any signal whose envelope is not constant. Backing off makes the
amplifier more linear but throws away power, and in a power-limited downlink every decibel of
OBO is a decibel of $C/N$. The optimum is where the \emph{total degradation}---the OBO plus the
extra $E_s/N_0$ the demodulator needs because of distortion---is smallest; we compute it for
DVB-S2 constellations in Section~\ref{sec:ch23:dvbs2}.

\subsection{Intermodulation in multicarrier transponders}

When one transponder carries many independent carriers (FDMA, Section~\ref{sec:ch23:access}), the
nonlinearity mixes them. A memoryless nonlinearity expanded in a power series has a cubic term
$a_3x^3$; with two carriers at $f_1$ and $f_2$ this produces \emph{third-order intermodulation
products} at $2f_1-f_2$ and $2f_2-f_1$, right next to the originals and inside the transponder.
With $N$ carriers the number of such products grows roughly as $N^3$, and they pile up in the
middle of the band like noise. Because the products scale as the cube of the carrier amplitude,
the carrier-to-intermodulation ratio $C/I_{\text{IM}}$ improves steadily with back-off,
approaching 2\,dB per decibel of back-off once the amplifier is well into its linear region. Figure~\ref{fig:ch23_intermod} shows a simulation: seven
QPSK carriers through a Saleh TWTA, with one channel slot left empty so the intermodulation that
falls into it can be measured.

\mdcfig{ch23_intermod}{Left: spectrum at the output of a TWTA carrying seven QPSK carriers at
three input back-offs; intermodulation fills the empty slot and spreads outside the band.
Right: carrier-to-intermodulation ratio in the empty slot and the resulting output back-off,
versus input back-off.}

At saturation the $C/I_{\text{IM}}$ is only about 11\,dB; at 6\,dB input back-off it is about
17\,dB, with the output about 3\,dB below saturation, and at 12\,dB it is 23\,dB with the output
7\,dB down. Multicarrier transponders therefore run backed off by
several decibels, and the intermodulation term appears explicitly in the link budget
(the 20\,dB used in the DTH example). A \emph{lineariser}---a predistorter ahead of the TWTA whose
characteristic is approximately the inverse of the tube's---gains back a decibel or two of
usable power, and is standard on modern payloads.

\subsection{Processing payloads}

A \emph{digital transparent processor} (DTP) digitises the uplink, channelises it with polyphase
filter banks (Chapter~\ref{ch:dsp}), routes any sub-channel from any input beam to any output
beam and frequency, and converts back to analog: still transparent to the waveform, but with the
connectivity of a switch. A \term{regenerative} payload goes further and demodulates and decodes
the uplink, switches packets and remodulates, gaining the error-probability combining of
Section~\ref{sec:ch23:combining}, the freedom to use different waveforms up and down, and the
ability to route traffic between satellites. Iridium was regenerative from the start, with
intersatellite links; Hughes' Spaceway~3 (2007) carried an on-board packet switch; and LEO
broadband constellations with laser crosslinks must process on board to route. The costs are
power, mass, radiation-hardened processing and the loss of future-proofing: a regenerative
satellite is tied to the waveform it was built for, unless it is software-defined. Recent
``software-defined'' satellites (Eutelsat Quantum, launched in 2021, was an early commercial
example) can reshape their beams, frequencies and power in orbit, blurring the line between
payload and network equipment.

\section{Multiple access}
\label{sec:ch23:access}

A satellite is a shared resource, and every multiple-access technique of
Chapter~\ref{ch:multipleaccess} has been tried on it. The special features of the satellite
channel---the long delay, the nonlinear transponder, and the huge disparity between a hub and
thousands of tiny terminals---shape the choice.

\paragraph{FDMA and SCPC.} In \term{frequency-division multiple access} each earth station
transmits its own carrier in its own slice of a transponder: \emph{single channel per carrier}
(SCPC) for one circuit, \emph{multiple channels per carrier} (MCPC) when a station multiplexes
several. FDMA is simple and needs no network-wide timing, but it forces the transponder to
back off to control intermodulation, and it is inflexible. It remains the standard for
point-to-point trunks, broadcast contribution links and the forward links of many networks.
Whether a multicarrier transponder is \emph{power-limited} or \emph{bandwidth-limited} is decided
by comparing the fraction of transponder power each carrier needs with the fraction of bandwidth
it occupies; a good plan uses both up at about the same time.

\paragraph{TDMA.} In \term{time-division multiple access} all stations share one wide carrier by
transmitting bursts in assigned time slots of a repeating frame. Only one carrier is present at a
time, so the transponder can run near saturation. Each burst starts with a preamble for carrier
and clock recovery and a unique word to mark its start, and bursts are separated by guard times;
the stations must control their transmit timing so that bursts arrive at the satellite in their
slots despite the satellite's motion within its station-keeping box, which changes the path by
up to about a hundred kilometres. Intelsat's TDMA system of the 1980s ran at about 120\,Mb/s.

\paragraph{MF-TDMA.} Modern VSAT return links combine the two: \term{multi-frequency TDMA} gives
the network a grid of carriers of different widths, each divided into time slots, and each
terminal hops from carrier to carrier according to a burst time plan broadcast by the hub
(Figure~\ref{fig:ch23_mftdma}). Capacity is assigned on demand (demand-assigned multiple access,
DAMA), so an idle terminal consumes almost nothing. The open standard is \term{DVB-RCS2}
(ETSI EN~301~545-2), which uses turbo-coded linear modulations from $\pi/2$-BPSK to 16QAM (and a
continuous-phase option for saturated amplifiers), with contention slots for random access and
logon.

\begin{figure}[htbp]\centering
\begin{tikzpicture}[x=0.52cm,y=0.5cm]
\foreach \c/\h/\y in {0/1/0,1/1/1,2/1/2,3/2/3,4/2/5}{
  \foreach \t in {0,...,19}{\draw[navy!30,fill=white] (\t,\y) rectangle ++(1,\h);}
}
% assigned bursts
\foreach \t/\y/\h/\col in {0/0/1/accent,1/0/1/accent,2/1/1/accent,5/2/1/accent,9/3/2/accent,
  3/0/1/practc,4/1/1/practc,6/0/1/practc,12/3/2/practc,13/3/2/practc,
  7/5/2/orange,8/5/2/orange,10/5/2/orange,15/1/1/orange,16/0/1/orange,
  11/2/1/purple,14/2/1/purple,17/2/1/purple,18/5/2/purple}{
  \fill[\col!55] (\t,\y) rectangle ++(1,\h);}
% contention slots
\foreach \t in {0,4,8,12,16}{\fill[pattern=north east lines,pattern color=gray] (\t,2) rectangle ++(1,1);}
\draw[navy,thick] (0,0) rectangle (20,7);
\draw[-{Stealth},thick] (0,-0.6) -- (20.8,-0.6) node[right,font=\scriptsize]{time};
\draw[-{Stealth},thick] (-0.6,0) -- (-0.6,7.6) node[above,font=\scriptsize]{frequency};
\node[font=\scriptsize,anchor=west] at (20.9,0.5) {narrow carriers};
\node[font=\scriptsize,anchor=west] at (20.9,6) {wide carriers};
\node[font=\scriptsize,anchor=west] at (20.9,2.5) {contention (hatched)};
\foreach \col/\lab/\x in {accent/terminal A/0,practc/terminal B/5,orange/terminal C/10,purple/terminal D/15}{
  \fill[\col!55] (\x,-2.0) rectangle ++(0.8,0.6); \node[font=\scriptsize,anchor=west] at (\x+0.9,-1.7) {\lab};}
\end{tikzpicture}
\caption{Multi-frequency TDMA on a VSAT return link. Carriers of different bandwidths are divided
into time slots; the hub's burst time plan tells each terminal which slot on which carrier to use in
every frame, so terminals hop in frequency and share capacity on demand. Some slots are reserved
for contention (random) access and logon.}
\label{fig:ch23_mftdma}
\end{figure}

\paragraph{CDMA and spreading.} Code-division multiple access (Chapter~\ref{ch:spreadspectrum})
lets carriers overlap in frequency and time. Globalstar's air interface was derived from IS-95; the
military uses spreading for jam resistance and low probability of intercept. A commercial use is
subtler: small aeronautical and mobile terminals \emph{spread} their return signals to lower their
power spectral density enough to meet the off-axis EIRP-density limits of
Section~\ref{sec:ch23:spectrum}, trading spectral efficiency for regulatory compliance.

\paragraph{Carrier-in-carrier.} A beautiful trick exploits the bent pipe. Two stations exchanging
traffic can transmit on the \emph{same} frequency at the same time; each receives the sum of its
partner's signal and its own, delayed by the round trip and distorted by the transponder. But each
station knows exactly what it sent, so it can estimate the delay, frequency offset and complex gain
of its own echo, synthesise it and subtract it, leaving the partner's signal. Commercial
implementations (``paired carrier multiple access'') nearly halve the transponder bandwidth needed
for a duplex link, at the cost of an adaptive canceller---the same idea as echo cancellation in
full-duplex modems (Chapter~\ref{ch:wireline}).

\paragraph{Random access: ALOHA and its descendants.} Many VSAT terminals send short, infrequent
messages---a card payment, a meter reading, a request for capacity. Assigning them slots in advance
wastes capacity; waiting half a second to ask for a slot wastes time. Norman Abramson's ALOHAnet
(Hawaii, 1971) proposed simply transmitting and retransmitting if a collision occurred. If packets
of duration $T$ arrive as a Poisson process with $G$ attempts per packet time, a packet succeeds
only if no other starts within $\pm T$ of it, so the throughput is
\begin{equation}
S=Ge^{-2G}\ \ \text{(pure ALOHA)},\qquad S=Ge^{-G}\ \ \text{(slotted ALOHA)},
\label{eq:ch23:aloha}
\end{equation}
with maxima $1/2e\approx18.4\%$ at $G=1/2$ and $1/e\approx36.8\%$ at $G=1$. ALOHA was a satellite
protocol before it was an Ethernet one. Its modern descendants break the $1/e$ barrier with
signal processing: in \emph{contention resolution diversity slotted ALOHA} (CRDSA), used in
DVB-RCS2, each packet is sent twice in random slots of a frame, with pointers to its twin; the
receiver decodes any clean copy, subtracts its replica from the twin's slot with successive
interference cancellation, and iterates, reaching peak throughputs above 0.5 packets per slot.
The same peeling idea as in the decoding of LDPC codes over the erasure channel
(Chapter~\ref{ch:moderncodes}) is at work.

\section{The DVB-S2 family}
\label{sec:ch23:dvbs2}

\subsection{From DVB-S to DVB-S2}

The first DVB satellite standard, \term{DVB-S} (ETSI EN~300~421), was designed in 1993--94
for digital television. It used QPSK with a roll-off of 0.35, and the concatenated coding that
deep-space missions had pioneered: an outer Reed--Solomon (204,188) code over bytes, a
convolutional interleaver, and an inner convolutional code of constraint length 7 punctured to
rates from 1/2 to 7/8, decoded with the Viterbi algorithm (Chapter~\ref{ch:classiccodes}). It
was robust, simple and enormously successful. But by 2003 turbo and LDPC codes had shown that the
Shannon limit was within reach (Chapter~\ref{ch:moderncodes}), and the DVB project set out to
design a successor with a new code, more efficient modulations and support for interactive
and professional services. The result, \term{DVB-S2}, was published as ETSI EN~302~307 in 2005
(now EN~302~307-1), and delivers roughly 30\% more capacity than DVB-S under the same conditions.
It is the physical layer of nearly all satellite TV launched since, of most VSAT forward links,
and of broadcast contribution links. Figure~\ref{fig:ch23_dvbs2_chain} shows the transmitter.

\begin{figure}[htbp]\centering
\begin{tikzpicture}[node distance=0.28cm and 0.28cm,
  blk/.style={draw=navy,fill=navy!6,thick,minimum height=1.25cm,text width=1.5cm,align=center,font=\sffamily\scriptsize},
  arr/.style={-{Stealth[length=2mm]},thick,navy}]
\node[blk] (ma) {mode adapt.\\CRC-8,\\BBHEADER};
\node[blk,right=of ma] (sa) {stream adapt.\\padding,\\scrambler};
\node[blk,right=of sa,fill=accent!8,draw=accent] (bch) {BCH outer\\encoder};
\node[blk,right=of bch,fill=accent!8,draw=accent] (ldpc) {LDPC inner\\encoder\\64800/16200};
\node[blk,right=of ldpc] (il) {bit\\interleaver};
\node[blk,below=0.75cm of il] (map) {mapper\\QPSK, 8PSK,\\16/32APSK};
\node[blk,left=of map] (pl) {PL framing:\\header, pilots,\\scrambler};
\node[blk,left=of pl] (rrc) {RRC filter\\$\alpha$ = 0.35,\\0.25, 0.20};
\node[blk,left=of rrc] (qm) {I/Q\\modulator\\to RF};
\node[left=0.45cm of ma,font=\sffamily\scriptsize,align=center] (in) {MPEG-TS\\or generic\\streams};
\draw[arr] (in) -- (ma); \draw[arr] (ma) -- (sa); \draw[arr] (sa) -- (bch); \draw[arr] (bch) -- (ldpc);
\draw[arr] (ldpc) -- (il); \draw[arr] (il) -- (map); \draw[arr] (map) -- (pl); \draw[arr] (pl) -- (rrc); \draw[arr] (rrc) -- (qm);
\node[below=0.08cm of bch,font=\sffamily\scriptsize,accent] {FEC: BCH + LDPC};
\node[above=0.08cm of ldpc,font=\sffamily\scriptsize,gray,align=center] {MODCOD chosen per frame (ACM)};
\end{tikzpicture}
\caption{The DVB-S2 transmitter (after ETSI EN~302~307-1). Input streams are packed into
baseband frames, protected by a BCH outer code and an LDPC inner code, interleaved, mapped to a
constellation, framed with a physical-layer header and optional pilots, and pulse-shaped.}
\label{fig:ch23_dvbs2_chain}
\end{figure}

\subsection{The pieces}

\paragraph{Framing.} Input packets are packed into \emph{baseband frames} (BBFRAMEs) with an
80-bit header describing the stream; each BBFRAME is exactly the payload of one FEC codeword.

\paragraph{FEC.} The inner code is an LDPC code with codeword length 64\,800 bits (normal frames)
or 16\,200 bits (short frames, for lower latency) and eleven code rates from 1/4 to 9/10. Its
parity-check matrix has an \emph{irregular repeat--accumulate} structure: the parity part is a
staircase, so encoding takes linear time, while the information part is built from cyclically
shifted blocks of 360 columns, which makes a partly parallel decoder straightforward to build in
hardware. Decoders run tens of iterations of belief propagation (Chapter~\ref{ch:moderncodes}). An
outer BCH code correcting 8 to 12 bit errors cleans up the residual errors that LDPC decoders
leave at high SNR (the ``error floor''), giving quasi-error-free output---in the standard's
definition, a packet error rate around $10^{-7}$.

\paragraph{Modulation.} QPSK and 8PSK with Gray mapping, and the ring constellations 16APSK
(4+12 points) and 32APSK (4+12+16), whose ring ratios are optimised for each code rate. The coded
bits of the higher-order constellations are interleaved so that the unequal protection of
different bit positions is spread over the LDPC code's variable nodes, a form of bit-interleaved
coded modulation.

\paragraph{Physical-layer framing.} Each codeword is mapped to an integer number of 90-symbol
\emph{slots} (a normal QPSK frame is 360 slots) and preceded by a 90-symbol \term{PLHEADER}: a
26-symbol start-of-frame sequence and a 64-symbol code carrying 7 bits of signalling---the
MODCOD, the frame length and whether pilots are present---protected by a (64,7) biorthogonal
Reed--Muller code and sent in robust $\pi/2$-BPSK. A receiver can therefore decode the header of
every frame, even one using a MODCOD it cannot demodulate, and so can find the frames meant for
it. Optional \emph{pilot blocks} of 36 known symbols after every 16 slots support carrier and
phase recovery at very low SNR and with high-order constellations (Chapter~\ref{ch:sync}).
Finally, everything after the header is multiplied by a complex scrambling sequence to whiten the
spectrum, and filtered with a root-raised-cosine pulse of roll-off 0.35, 0.25 or 0.20
(Chapter~\ref{ch:baseband}).

\begin{table}[htbp]\centering\small
\caption{The 28 DVB-S2 MODCODs: spectral efficiency (information bits per transmitted symbol,
normal frames, no pilots) and ideal $E_s/N_0$ for quasi-error-free operation in AWGN, after ETSI
EN~302~307-1. Real receivers need a few tenths of a decibel more.}\label{tab:ch23:dvbs2}
\begin{tabular}{@{}l r r @{\hspace{2.2em}} l r r@{}}
\toprule
MODCOD & b/sym & $E_s/N_0$ (dB) & MODCOD & b/sym & $E_s/N_0$ (dB)\\
\midrule
QPSK 1/4 & 0.490 & $-2.35$ & 8PSK 3/5 & 1.779 & 5.50\\
QPSK 1/3 & 0.656 & $-1.24$ & 8PSK 2/3 & 1.980 & 6.62\\
QPSK 2/5 & 0.789 & $-0.30$ & 8PSK 3/4 & 2.228 & 7.91\\
QPSK 1/2 & 0.989 & 1.00 & 8PSK 5/6 & 2.479 & 9.35\\
QPSK 3/5 & 1.188 & 2.23 & 8PSK 8/9 & 2.646 & 10.69\\
QPSK 2/3 & 1.322 & 3.10 & 8PSK 9/10 & 2.679 & 10.98\\
QPSK 3/4 & 1.487 & 4.03 & 16APSK 2/3 & 2.637 & 8.97\\
QPSK 4/5 & 1.587 & 4.68 & 16APSK 3/4 & 2.967 & 10.21\\
QPSK 5/6 & 1.655 & 5.18 & 16APSK 4/5 & 3.166 & 11.03\\
QPSK 8/9 & 1.767 & 6.20 & 16APSK 5/6 & 3.300 & 11.61\\
QPSK 9/10 & 1.789 & 6.42 & 16APSK 8/9 & 3.523 & 12.89\\
32APSK 3/4 & 3.703 & 12.73 & 16APSK 9/10 & 3.567 & 13.13\\
32APSK 4/5 & 3.952 & 13.64 & 32APSK 8/9 & 4.398 & 15.69\\
32APSK 5/6 & 4.120 & 14.28 & 32APSK 9/10 & 4.453 & 16.05\\
\bottomrule
\end{tabular}
\end{table}

Table~\ref{tab:ch23:dvbs2} and Figure~\ref{fig:ch23_dvbs2_modcods} show the resulting menu. The
MODCODs span 18\,dB of $E_s/N_0$ in steps of typically 0.5--1\,dB, with efficiencies from 0.49 to
4.45\,b/s/Hz. At low rates they sit about 1\,dB from the Shannon bound $\log_2(1+E_s/N_0)$; the
gap grows to nearly 3\,dB for 32APSK 9/10, the cost of a finite code length, a quasi-error-free
target, and constellations that are not shaped for capacity.

\mdcfig{ch23_dvbs2_modcods}{The DVB-S2 MODCODs against the Shannon bound. The staircase is the
efficiency an adaptive (ACM) link obtains by always using the most efficient MODCOD that closes
at the current $E_s/N_0$ (\texttt{satellite.acm\_select}).}

\subsection{Adaptive coding and modulation}

DVB-S2 allows three modes. \emph{Constant} coding and modulation (CCM) uses one MODCOD for
everything, as broadcasting must, since every receiver in the footprint gets the same signal.
\emph{Variable} coding and modulation (VCM) assigns different MODCODs to different services,
for example a robust MODCOD for standard-definition and a more efficient one for
high-definition streams. \textbf{Adaptive coding and modulation}\index{adaptive coding and modulation} (ACM) changes the MODCOD frame by
frame according to the reported conditions of the receiver each frame is addressed to. Because
every PLHEADER announces its MODCOD, a receiver simply skips frames it cannot or need not
decode. In an interactive network each terminal measures its $E_s/N_0$ from the headers and
pilots and reports it over the return link; the gateway's scheduler chooses the MODCOD for that
terminal's traffic, with a hysteresis margin that covers the fade that can develop during the
report-and-react loop---about half a second over GEO. The Ka-band worked example showed ACM
multiplying the average capacity by 3.7 at unchanged availability, and it also allows a
network to serve a mix of small and large terminals, and terminals at beam centre and beam
edge, each at its best rate.

\begin{keyidea}[Match the rate to the channel, not to the worst case]
Fixed links must be designed for the worst few hours of the year and waste their margin the rest
of the time. ACM turns margin into throughput whenever the sky is clear, and gives it back only
when needed. The same principle runs through the whole book: rate adaptation in Wi-Fi and
cellular (Chapters~\ref{ch:wifi} and~\ref{ch:cellular}), bit loading in DSL and OFDM
(Chapter~\ref{ch:ofdm}), water-filling in information theory (Chapter~\ref{ch:infotheory}).
\end{keyidea}

\subsection{Why APSK?}

Why did DVB-S2 choose rings of points rather than the square QAM grids used almost everywhere
else? Because of the amplifier. A satellite downlink is power-limited, so the TWTA should run as
close to saturation as possible, where its AM/AM compression and AM/PM rotation distort any
constellation with several amplitude levels. 16APSK has only two rings, with most of its points on
the outer one: its peak-to-average power ratio at the symbol instants is 1.1\,dB against 2.6\,dB
for 16QAM, and each ring is compressed and rotated as a whole, so the transmitter can predistort
simply by adjusting two radii and one relative phase, and the receiver can track the distorted
ring positions with its pilots. Figure~\ref{fig:ch23_apsk_vs_qam} makes the comparison with a
simulation: pulse-shaped 16QAM and 16APSK through the Saleh TWTA, with the \emph{total degradation}
(output back-off plus the extra $E_s/N_0$ needed, treating the distortion as additional Gaussian
noise, for a rate-3/4 code whose threshold is 10.2\,dB) plotted against output back-off.

\mdcfig{ch23_apsk_vs_qam}{16QAM versus DVB-S2 16APSK through a TWTA. Left and centre: symbol-level
constellations before (open) and after (filled) the amplifier driven at 0\,dB input back-off.
Right: total degradation versus output back-off for RRC-filtered signals (roll-off 0.2), for a
receiver that corrects only gain and phase (solid) and one that learns the distorted point
positions (dashed).}

The results are instructive. Every curve has a minimum: too little back-off and distortion
dominates, too much and power is wasted. With a receiver that corrects only a common gain and
phase, 16APSK's optimum total degradation is about 0.5\,dB lower than 16QAM's. A receiver that
learns the actual (distorted) centroid of every point closes most of that gap, and both reach a
total degradation of about 1.6\,dB at 1\,dB of output back-off. The remaining degradation comes
from the \emph{memory} that pulse shaping creates---the amplifier sees the filtered waveform, not
the symbols---which symbol-level fixes cannot remove. The lesson is that APSK's advantage is real
but moderate for 16 points, and that it depends on how clever the ground equipment is; its
greater benefit is at higher orders and in making simple predistortion possible, and DVB-S2X
extended the family to 64-, 128- and 256-point APSK.

\subsection{DVB-S2X and the return link}

The 2014 extension \term{DVB-S2X} (ETSI EN~302~307-2) added sharper roll-offs of 0.15, 0.10 and
0.05---a useful few per cent of capacity on a clean channel, since the occupied bandwidth is
$R_s(1+\alpha)$---many more MODCODs with finer steps, constellations up to 256APSK for
high-SNR professional links, \emph{very-low-SNR} MODCODs reaching down to about $-10$\,dB $E_s/N_0$
for small mobile antennas, and a \emph{superframe} structure designed to support beam hopping and
multibeam precoding. For the return link from interactive terminals, \term{DVB-RCS2}
(ETSI EN~301~545-2, 2012) specifies MF-TDMA with turbo-coded bursts, as in
Section~\ref{sec:ch23:access}. Many commercial VSAT systems use proprietary variants of the same
ideas; the CCSDS standards (Section~\ref{sec:ch23:deepspace}) govern space agencies' links.

\section{Satellite systems today}
\label{sec:ch23:systems}

\subsection{VSAT networks}

\begin{figure}[htbp]\centering
\begin{tikzpicture}[>={Stealth[length=2.2mm]},
  term/.style={draw=navy,fill=navy!6,thick,rounded corners=2pt,minimum width=1.3cm,minimum height=0.65cm,align=center,font=\sffamily\scriptsize}]
\node[draw=navy,fill=navy!12,thick,minimum width=1.2cm,minimum height=0.5cm,font=\sffamily\scriptsize] (sat) at (0,3.3) {satellite};
\draw[navy,thick] (sat.west) -- ++(-0.6,0.25) -- ++(0,-0.5) -- cycle;
\draw[navy,thick] (sat.east) -- ++(0.6,0.25) -- ++(0,-0.5) -- cycle;
\node[term,minimum width=2.1cm,minimum height=0.9cm] (hub) at (-4.2,0) {hub / gateway\\(large antenna)};
\node[term] (v1) at (-0.9,-0.2) {VSAT\\shop};
\node[term] (v2) at (1.1,-0.4) {VSAT\\cell site};
\node[term] (v3) at (3.1,-0.2) {VSAT\\ship};
\node[term] (v4) at (5.0,0.1) {VSAT\\school};
\node[draw=gray,thick,rounded corners=2pt,font=\sffamily\scriptsize,align=center] (net) at (-6.6,0) {terrestrial\\network,\\Internet};
\draw[thick,gray] (net) -- (hub);
\draw[->,thick,accent] (hub.north) to[bend left=10] node[left,font=\sffamily\scriptsize,align=right,pos=0.45]{forward: one\\DVB-S2 TDM\\carrier (ACM)} (sat.south west);
\draw[->,thick,practc] (sat.south west) ++(0.25,0) to[bend left=10] ($(hub.north)+(0.45,0)$);
\foreach \v in {v1,v2,v3,v4}{
  \draw[->,accent] ($(sat.south)+(0.1,0)$) -- (\v.north);
}
\draw[->,practc,dashed] (v4.north west) to[bend left=12] ($(sat.south east)+(-0.1,0)$);
\node[font=\sffamily\scriptsize,practc,align=left] at (4.4,2.1) {return: MF-TDMA\\bursts (DVB-RCS2),\\on demand};
\end{tikzpicture}
\caption{A star VSAT network. The hub transmits one wide TDM carrier (DVB-S2 with ACM) that every
terminal receives, picking out its own frames; terminals share MF-TDMA return carriers assigned
by the hub. Terminal-to-terminal traffic passes through the hub (a double hop).}
\label{fig:ch23_vsat}
\end{figure}

A \term{VSAT} (very small aperture terminal) network links hundreds to hundreds of thousands of
small terminals---typically 0.75--2.4\,m dishes---through a satellite to a central \emph{hub}
(Figure~\ref{fig:ch23_vsat}). The architecture is asymmetric by design: the hub has a big antenna
and a powerful amplifier, so the \emph{forward link} can be one wide carrier that saturates a
transponder; the terminals are cheap and weak, so the \emph{return link} uses many narrow
MF-TDMA carriers. Terminal-to-terminal traffic in a star network makes a double hop through the
hub, over a second of round trip at GEO; \emph{mesh} networks let terminals with enough $G/T$ and
EIRP talk directly. VSATs carry lottery and card transactions for retailers and banks,
connect oil fields, mines and ships, provide rural schools and clinics with Internet, and---a large
market---\emph{backhaul} cellular base stations in places where fibre and microwave cannot
reach.

\begin{inpractice}[Cellular backhaul by satellite]
A rural 4G base station may carry only tens of megabits per second at its busiest hour, but it must
be available all the time, and many such sites are hundreds of kilometres from fibre. Satellite
backhaul, historically over GEO and increasingly over MEO and LEO, connects them. The engineering
is a blend of this chapter and Chapter~\ref{ch:cellular}: header compression and traffic
optimisation to save expensive satellite capacity, careful handling of the base station's timing
(which must come from GNSS rather than from the backhaul, since packet delay over the satellite
varies), and the delay budget of the cellular core, which tolerates the extra half-second of a
GEO link for data but makes voice calls between two satellite-backhauled sites painfully slow.
\end{inpractice}

\subsection{Ships and aircraft}

For decades Inmarsat's L-band services provided narrowband voice and data at sea and in the air,
with small, nearly omnidirectional antennas. Broadband mobility moved to Ku- and Ka-band in the
2010s (Inmarsat's Global Xpress Ka-band system began service in 2015), which requires directional
antennas that stay locked on a satellite while the platform rolls, pitches and turns: gyroscopically
stabilised dishes under radomes on ships, and low-profile mechanically or electronically steered
arrays on aircraft fuselages, where drag matters and the satellite may be only a few degrees above
the wing at high latitudes. Regulators treat these as \emph{earth stations in motion} and require
them to keep their off-axis emissions within limits despite motion, so their pointing systems must
mute the transmitter within a fraction of a second if pointing is lost. Aircraft speeds add a
Doppler of only about 10\,kHz at 12\,GHz---small compared to LEO---but LEO constellations have
now entered this market as well, with phased-array terminals on airliners and ships.

\subsection{GEO high-throughput satellites}

The GEO industry's answer to LEO is ever larger spot-beam satellites. Hughes' Jupiter~3
(EchoStar~XXIV, launched 2023) was designed for about 500\,Gb/s; Viasat's ViaSat-3 series was
designed for on the order of 1\,Tb/s per satellite with three satellites covering most of the
world, although the first, launched in 2023, suffered a reflector deployment failure that sharply
reduced its capacity---a reminder that everything on a satellite must work first time. SES's
O3b mPOWER in MEO uses steerable beams and digital processing to deliver capacity where it is
sold. The economics are governed by cost per delivered bit, and the capacity arithmetic of
Section~\ref{sec:ch23:antennas}: more beams, more reuse, more feeder-link spectrum and ACM.

\subsection{LEO mega-constellations}

\begin{figure}[htbp]\centering
\begin{tikzpicture}[>={Stealth[length=2.2mm]},
  sat/.style={draw=navy,fill=navy!12,thick,minimum width=0.9cm,minimum height=0.45cm,font=\sffamily\scriptsize},
  gnd/.style={draw=navy,fill=navy!5,thick,rounded corners=2pt,align=center,font=\sffamily\scriptsize,minimum height=0.7cm}]
\draw[navy!30,thick] (-6.3,-0.25) -- (6.3,-0.25);
\node[font=\sffamily\scriptsize,navy!60] at (5.6,-0.55) {Earth surface};
\node[sat] (s1) at (-4,2.6) {sat A};
\node[sat] (s2) at (-1.2,3.0) {sat B};
\node[sat] (s3) at (1.6,3.0) {sat C};
\node[sat] (s4) at (4.3,2.6) {sat D};
\draw[<->,thick,accent,densely dashed] (s1) -- node[above,font=\sffamily\scriptsize]{laser ISL} (s2);
\draw[<->,thick,accent,densely dashed] (s2) -- (s3);
\draw[<->,thick,accent,densely dashed] (s3) -- node[above,font=\sffamily\scriptsize]{laser ISL} (s4);
\node[gnd] (ut) at (-4.6,0.3) {user terminal\\(phased array)};
\node[gnd] (ut2) at (-2.2,0.3) {user terminal};
\node[gnd] (gw) at (4.6,0.3) {gateway\\(Ka/E-band)};
\node[gnd] (pop) at (4.6,-1.3) {point of presence\\$\to$ Internet};
\draw[<->,thick,navy] (ut) -- node[left,font=\sffamily\scriptsize]{Ku} (s1);
\draw[<->,thick,navy] (ut2) -- (s1);
\draw[<->,thick,navy] (gw) -- node[right,font=\sffamily\scriptsize]{Ka} (s4);
\draw[thick,gray] (gw) -- (pop);
\draw[->,thick,practc,dotted] (-4.0,1.9) -- node[below,font=\sffamily\scriptsize,sloped]{handover} (-1.4,2.4);
\end{tikzpicture}
\caption{A LEO broadband constellation with intersatellite links. Traffic from a user terminal goes
up to the satellite in view, is routed across laser crosslinks, and comes down to a gateway (or,
over oceans, to a gateway thousands of kilometres away). Users hand over between satellites every
few minutes or less.}
\label{fig:ch23_leo_arch}
\end{figure}

The new LEO systems combine everything in this chapter. They are built as \emph{Walker}
constellations: several \emph{shells}, each a set of orbital planes at a common altitude and
inclination with satellites evenly spaced in each plane. Inclined shells (around 53$^\circ$)
concentrate satellites over the mid-latitudes where most customers live; polar shells cover
the high latitudes. The coverage calculation of Section~\ref{sec:ch23:geometry} says that a few
hundred satellites suffice to have one in view almost everywhere, but capacity, not coverage,
drives the numbers: each satellite has a fixed budget of power, beams and spectrum, and one
beam serves one ``cell'' on the ground, so capacity per square kilometre is limited in dense
areas, while most satellites spend most of their time over oceans and deserts with little to do.

\paragraph{Links.} User links are in Ku-band (for Starlink, downlink 10.7--12.7\,GHz, uplink
14.0--14.5\,GHz) with phased arrays on both ends; gateways use Ka-band and, increasingly, E-band.
OneWeb's first generation is a bent pipe with 16 fixed, elongated user beams per satellite, so a
gateway must be in view of every satellite that carries traffic. Starlink's satellites process on
board and are connected by \term{intersatellite laser links} (ISLs), each reportedly running at
tens of gigabits per second or more, so traffic can be routed through space to a distant gateway:
coverage of oceans, polar regions and places without gateways becomes possible, and for very long
routes the path through vacuum can even beat fibre, in which light travels at only about two
thirds of its vacuum speed.

\paragraph{Mobility management.} The network must track which satellite, beam and frequency serves
each terminal, reassign them as satellites move, and route packets through a topology that changes
every few seconds. Measurement studies of Starlink by academic researchers have observed periodic
latency changes on a timescale of about 15 seconds, consistent with a global scheduler that
reassigns terminals at fixed intervals. Handovers are planned in advance from known orbits, so a
terminal can steer its beam to the next satellite at the right instant; the electronic beam
steering of Section~\ref{sec:ch23:antennas} is what makes this possible.

\paragraph{Sharing the sky.} A LEO system must protect GEO networks (the EPFD limits and GSO arc
avoidance of Section~\ref{sec:ch23:spectrum}), coordinate with other NGSO systems, limit
interference to radio astronomy, and dispose of its satellites safely. These constraints, as much
as physics, shape where beams can point and how much capacity reaches the ground near the
equator.

\section{5G from space: non-terrestrial networks}
\label{sec:ch23:ntn}

For fifty years satellite and cellular systems were designed by different communities with
different standards. 3GPP's \term{non-terrestrial network} (NTN) work, a study from Release~15 and
normative from \textbf{Release~17} (completed in 2022), changed that by adapting the 5G NR and
IoT air interfaces (Chapter~\ref{ch:cellular}) to satellite links. Release~17 introduced
NR-NTN in two FR1 bands, n255 in the L-band and n256 in the S-band, and IoT-NTN (NB-IoT and
LTE-M over satellite), for transparent payloads in LEO and GEO. Later releases added Ka-band
operation for VSAT-class terminals, coverage and mobility enhancements, network verification of
the terminal's reported location (Release~18) and regenerative payloads with the base station on
board (Release~19).

\begin{figure}[htbp]\centering
\begin{tikzpicture}[node distance=0.35cm and 0.5cm,
  blk/.style={draw=navy,fill=navy!6,thick,rounded corners=2pt,minimum height=0.7cm,minimum width=1.3cm,align=center,font=\sffamily\scriptsize},
  sat/.style={draw=accent,fill=accent!8,thick,minimum height=0.7cm,minimum width=1.3cm,align=center,font=\sffamily\scriptsize},
  arr/.style={<->,thick,navy}]
\node[blk] (ue1) {UE};
\node[sat,right=1.5cm of ue1] (s1) {satellite\\(transparent)};
\node[blk,right=1.5cm of s1] (gw1) {NTN\\gateway};
\node[blk,right=of gw1] (gnb1) {gNB};
\node[blk,right=of gnb1] (cn1) {5G core};
\draw[arr] (ue1) -- node[above,font=\sffamily\scriptsize]{service} node[below,font=\sffamily\scriptsize]{link} (s1);
\draw[arr] (s1) -- node[above,font=\sffamily\scriptsize]{feeder} node[below,font=\sffamily\scriptsize]{link} (gw1);
\draw[arr] (gw1) -- (gnb1); \draw[arr] (gnb1) -- (cn1);
\node[blk,below=0.9cm of ue1] (ue2) {UE};
\node[sat,right=1.5cm of ue2] (s2) {satellite\\with gNB};
\node[blk,right=1.5cm of s2] (gw2) {NTN\\gateway};
\node[blk] (cn2) at (cn1 |- gw2) {5G core};
\draw[arr] (ue2) -- node[above,font=\sffamily\scriptsize]{NR Uu} (s2);
\draw[arr] (s2) -- node[above,font=\sffamily\scriptsize]{NG/SRI} (gw2);
\draw[arr] (gw2) -- (cn2);
\node[left=0.2cm of ue1,font=\sffamily\scriptsize,align=right] {Rel-17\\transparent};
\node[left=0.2cm of ue2,font=\sffamily\scriptsize,align=right] {Rel-19\\regenerative};
\end{tikzpicture}
\caption{3GPP NTN architectures. In the transparent case (Release~17) the satellite is an RF repeater
and the whole NR protocol stack runs on the ground, so the air-interface delay includes the feeder
link. In the regenerative case (Release~19) the gNB is on board and the feeder link carries the
interface to the core network over a satellite radio interface (SRI).}
\label{fig:ch23_ntn}
\end{figure}

\subsection{What had to change}

NR was designed for cells a few kilometres across, where delays are microseconds and Doppler is
set by car and train speeds. Three things break over a satellite.

\paragraph{Timing.} An NR uplink must arrive at the base station aligned with its slot structure,
so each terminal transmits early by a \emph{timing advance} equal to its round-trip delay. On the
ground this is at most a millisecond or so; through a transparent LEO satellite at 600\,km the
round trip reaches about 26\,ms at 10$^\circ$ elevation, and through GEO about 541\,ms (TR~38.821),
far beyond the range of the terrestrial mechanism. NTN terminals therefore compute their own
timing advance: every NTN terminal has a GNSS receiver, and the cell broadcasts the satellite's
ephemeris and a \emph{common} timing advance for the feeder-link part of the path (in system
information block SIB19), so the terminal adds its own service-link delay from geometry and
closes the loop with small corrections from the network. Timing relationships that assume a
reply within a few slots---``transmit PUSCH $K_2$ slots after the grant''---gain a
scheduling offset $K_{\text{offset}}$ covering the round trip.

\paragraph{Frequency.} From Section~\ref{sec:ch23:doppler}, a 600\,km satellite at 2\,GHz produces
up to about $\pm46$\,kHz of Doppler (23\,ppm, consistent with the figures in TR~38.811) changing by
up to about 600\,Hz/s---three 15\,kHz subcarriers of offset. The terminal \emph{pre-compensates}
its uplink frequency using its known position and velocity and the satellite ephemeris, so the
signal arrives at the satellite nearly on frequency; on the downlink, the network can remove the
Doppler common to the beam centre, and the terminal tracks the residual.

\paragraph{HARQ.} NR's hybrid ARQ uses stop-and-wait processes: the transmitter must wait a round
trip for each acknowledgement, and it keeps the channel busy by running several processes in
parallel. The number needed is roughly the round trip divided by the slot duration---hundreds over
GEO. Release~17 raised the maximum number of HARQ processes to 32 and allowed the network to
\emph{disable HARQ feedback} per process, relying instead on blind repetitions, robust MODCODs and
the retransmissions of higher layers (RLC ARQ), whose timers were lengthened accordingly.

\begin{worked}[Timing and Doppler for an S-band NTN cell]
A transparent LEO satellite at 600\,km serves a beam in which terminals see it at elevations
between 30$^\circ$ and 90$^\circ$. From~\eqref{eq:ch23:slant}, the service-link range varies from
600 to 1075\,km, a one-way delay difference of $475/299\,792\approx1.6$\,ms, and the round-trip
difference within the beam is 3.2\,ms---more than three NR slots at 15\,kHz. A single common
timing advance cannot serve the whole beam, which is why each terminal computes its own from GNSS
and ephemeris. At 2\,GHz the Doppler at the beam edges differs by tens of kilohertz; if the network
pre-compensated only the beam-centre Doppler, terminals at the edge would still be several
subcarriers off. Each terminal therefore corrects its own. With a GNSS position error of 10\,m and
an ephemeris error of 100\,m, the residual timing error is well under a microsecond, a small
fraction of the 4.7\,$\mu$s cyclic prefix, and the residual Doppler error from a velocity error of
0.1\,m/s is under 1\,Hz.
\end{worked}

\subsection{Direct to device}

The most visible new satellite service is \term{direct-to-device} (D2D): links from satellites to
ordinary phones. There are two approaches. Standards-based NTN handsets implement the Release~17
features above in satellite bands such as n255 and n256. The alternative, used by the
supplemental-coverage services that US regulators authorised from 2024, reuses a mobile
operator's terrestrial spectrum and serves \emph{unmodified} phones: the satellite carries a
base station that makes itself look like a distant terrestrial cell, pre-compensating the
Doppler and delay to the centre of each beam so that the phone's standard procedures work.

The link budget explains why D2D satellites are so large.

\begin{worked}[Why direct-to-cell satellites have huge antennas]
Take a 600\,km satellite serving a phone at 30$^\circ$ elevation (range 1075\,km) at 2\,GHz:
$L_{\text{fs}}=159.1$\,dB. Allow 3\,dB for the polarisation mismatch between a linearly polarised
phone antenna and a circularly polarised satellite, and 1\,dB for other losses. The phone's antenna
and the user's body give about $-3$\,dBi, and with a 5\,dB noise figure and 290\,K antenna
temperature the phone's $T_{\text{sys}}\approx917$\,K, so its $G/T\approx-32.6$\,dB/K---some
50\,dB worse than a Ku-band dish.

\emph{Downlink.} For a 5\,MHz LTE carrier at 5\,dB SNR the phone needs
$C/N_0=67+5=72$\,dB-Hz. From~\eqref{eq:ch23:cn0}, the satellite must radiate
$\text{EIRP}=72+159.1+4+32.6-228.6=39.1$\,dBW per beam. A 10\,m$^2$ array at 2\,GHz has a gain of
about 35.9\,dBi, so about 2\,W of RF per beam suffices. The downlink closes comfortably.

\emph{Uplink.} The phone radiates 23\,dBm with $-3$\,dBi: an EIRP of $-10$\,dBW, fixed by
regulation and battery. With the 10\,m$^2$ array ($G/T\approx8.9$\,dB/K for a 500\,K
system temperature looking at the warm Earth),
$C/N_0=-10-159.1-4+8.9+228.6=64.4$\,dB-Hz: about 12\,dB of SNR in one 180\,kHz resource block, but
$-2.6$\,dB across 5\,MHz. With a 1\,m$^2$ array the uplink has only 54\,dB-Hz, enough for text
messages at low rate. The uplink, not the downlink, sets the aperture; and because each beam's
footprint shrinks as the aperture grows, large arrays also bring the frequency reuse that
capacity requires. Hence satellites with arrays of tens of square metres, and messaging-only
services from smaller ones.
\end{worked}

\section{Deep-space communication}
\label{sec:ch23:deepspace}

At the far end of the scale are the links to planetary spacecraft, where distances are measured in
astronomical units ($1\,\text{AU}=1.496\times10^8$\,km) and received powers in attoWatts. NASA's
\term{Deep Space Network} (DSN) has three complexes, near Goldstone (California), Madrid and
Canberra, roughly 120$^\circ$ apart in longitude so that a spacecraft is always in view of one of
them. Each has a 70\,m antenna and several 34\,m antennas, with cryogenically cooled low-noise
amplifiers giving system noise temperatures of a few tens of kelvin at X-band (downlink near
8.4\,GHz) and Ka-band (near 32\,GHz). The 70\,m antenna at Canberra is the only one able to send
commands to Voyager~2, which is far below the plane of the planets and visible only from the
southern hemisphere. Antennas can be \emph{arrayed}, combining their signals coherently to gain up
to $10\log_{10}N$\,dB. The same links provide navigation: two-way coherent Doppler measures a
spacecraft's line-of-sight velocity to a fraction of a millimetre per second, and ranging codes
measure its distance to metres---the synchronisation techniques of Chapter~\ref{ch:sync}
applied across the solar system.

\begin{history}[Coding theory's proving ground]
Deep space was the first customer for powerful error-correcting codes, because there every decibel
of coding gain translated directly into data rate. The Mariner missions to Mars from 1969 sent
pictures protected by a (32,6) biorthogonal Reed--Muller code, decoded on the ground by a fast
Hadamard transform in a machine known as the ``Green Machine''. Pioneer missions used long convolutional codes with
sequential decoding. The Voyagers used a constraint-length-7, rate-1/2 convolutional code with
Viterbi decoding, later concatenated with a Reed--Solomon (255,223) code---the pair that became the
DVB-S and deep-space standard for two decades. When Galileo's high-gain antenna failed to unfurl in
1991, engineers rescued the mission with its low-gain antenna, new on-board data compression and
more powerful concatenated codes, recovering most of the science at a data rate roughly a thousand
times lower than planned. The CCSDS (Consultative Committee for Space Data Systems) standardised
turbo codes shortly after their invention, and later a family of LDPC codes (CCSDS~131.0-B), now used
by missions throughout the solar system.
\end{history}

\begin{worked}[Voyager~1 at one light-day]
In late 2026 Voyager~1 will be roughly 173\,AU, or $2.59\times10^{10}$\,km, from Earth: light will take
about 24 hours to make the trip, and NASA expects the spacecraft to pass one light-day from Earth
around November 2026. Its X-band transmitter produces about 20\,W (13\,dBW) at 8.42\,GHz into a
3.7\,m dish of perhaps 55\% efficiency (47.7\,dBi): an EIRP of about 60.7\,dBW. The free-space loss is
\[
L_{\text{fs}}=20\log_{10}\frac{4\pi\cdot2.59\times10^{13}\cdot8.42\times10^9}{2.998\times10^8}=319.2\ \text{dB}.
\]
A 70\,m DSN antenna at 70\% efficiency has 74.3\,dBi; with $T_{\text{sys}}=20$\,K its $G/T$ is
61.3\,dB/K. So
\[
C/N_0=60.7-319.2+61.3+228.6=31.4\ \text{dB-Hz},
\]
and at the 160\,b/s that Voyager now sends, $E_b/N_0=31.4-22.0=9.4$\,dB before losses. Real links
lose several decibels to the residual carrier that the receiver tracks, to pointing, polarisation
and circuit losses and to weather; allowing about 4\,dB leaves roughly 5.4\,dB, against the
2--3\,dB that the concatenated code needs. A 34\,m antenna, 7\,dB less sensitive, could not close
the link alone---which is why Voyager's science data come through the 70\,m dishes or arrays of
34\,m ones. Figure~\ref{fig:ch23_deep_space} extends the calculation to the whole solar system.
Received power here is a few times $10^{-19}$\,W.
\end{worked}

\mdcfig{ch23_deep_space}{Supportable bit rate versus distance for a Voyager-like X-band spacecraft
(20\,W, 3.7\,m dish) received by a 70\,m or a 34\,m DSN antenna, assuming 2.5\,dB required
$E_b/N_0$ and 7\,dB for losses and margin. The rate falls as $1/d^2$: 20\,dB per decade of distance.
The curves are a model with rounded parameters, not a mission's actual link design.}

The $1/d^2$ law of Figure~\ref{fig:ch23_deep_space} is unforgiving, and the only escapes are
higher frequencies (a fixed dish has more gain at shorter wavelengths, so moving from X- to Ka-band
gains roughly 11\,dB of antenna gain at each end, less weather losses), bigger or arrayed antennas,
more power, and better codes. The biggest step is to optical wavelengths. NASA's \term{Deep Space
Optical Communications} (DSOC) experiment, launched in October 2023 on the Psyche asteroid mission,
sent data from a 22\,cm telescope with a near-infrared laser to the 5\,m Hale telescope at Palomar,
guided by an uplink laser beacon, and demonstrated rates up to about 267\,Mb/s from some 30 million
kilometres---tens of times what radio systems of similar size and power achieve at such
distances---and continued to link from beyond the orbit of Mars. The receiver counts individual
photons with superconducting nanowire detectors, and the modulation is pulse-position
modulation, the power-efficient choice when photons are few (Chapter~\ref{ch:infotheory}).

\section{Navigation satellites}

The most widely used satellite service is not communication at all but navigation. The GPS
constellation of about 30 satellites in six planes of 20\,200\,km MEO orbits, and its
counterparts Galileo, GLONASS and BeiDou, broadcast precisely timed spread-spectrum signals from
which a receiver solves for its position and time. GPS signals arrive roughly 19\,dB below the
thermal noise in their 2\,MHz bandwidth and are recovered by correlation, which makes them the ideal
case study for spread spectrum; they are treated in Chapter~\ref{ch:spreadspectrum}. They also
matter to every system in this chapter: NTN terminals compute timing advance and Doppler
pre-compensation from GNSS, cellular base stations and VSAT hubs take their time from it, and a
LEO user terminal uses it to know where to point.

\section{Looking back, looking up}

A satellite link is a long, clean, noisy, nonlinear, shared, rain-soaked pipe with a delay you can
hear. The engineering of it touches every chapter of this book: noise temperature and link
budgets (Chapter~\ref{ch:noise}), pulse shaping and matched filtering (Chapter~\ref{ch:baseband}),
constellations chosen for amplifiers (Chapter~\ref{ch:modulation}), synchronisation at very low SNR
(Chapter~\ref{ch:sync}), codes that approach capacity (Chapter~\ref{ch:moderncodes}), multibeam
antennas and MIMO precoding (Chapter~\ref{ch:mimo}), multiple access
(Chapter~\ref{ch:multipleaccess}) and, now, 5G itself (Chapter~\ref{ch:cellular}). The industry
has been declared mature, or dying, several times; each time new engineering has changed the
trade between altitude, antenna and cost, and a new generation of systems has appeared. The
integration of satellites into the terrestrial network, so that a phone moves seamlessly between a
tower and a satellite, is one of the central themes of the road to 6G.

\begin{labbox}[Satellite link budgets, geometry and DVB-S2]
The lab \texttt{moderncomms-labs/labs/lab18\_satellite\_link.py} builds on
\texttt{commlib/satellite.py}, the module that drew every data figure in this chapter. It
contains an interactive link-budget calculator (EIRP, path loss, rain, $G/T$, interference
combining) that reproduces the DTH and Ka-band worked examples; a GEO-versus-LEO geometry
explorer with look angles, pass duration, delay and Doppler; a rain-fade availability study using
the P.618 procedure; an ACM simulation that picks DVB-S2 MODCODs from a year of synthetic rain and
compares throughput with CCM; and a TWTA experiment in which you drive 16QAM and 16APSK through
the Saleh model, measure the total degradation and try your own predistorter. Start by
reproducing Figure~\ref{fig:ch23_leo_doppler}, then change the altitude and frequency to those of a
direct-to-cell system.
\end{labbox}

\problems
\begin{enumerate}[label=\thechapter.\arabic*]
\item Derive the orbital period of a circular orbit from Newton's law of gravitation, and compute
the altitude of a 12-hour circular orbit and of a circular orbit whose ground track repeats after
exactly 15 orbits per sidereal day (ignore $J_2$). Why are repeating ground tracks convenient for
LEO system planning?
\item Using the triangle of Figure~\ref{fig:ch23_geometry_tikz}, derive~\eqref{eq:ch23:slant}
and~\eqref{eq:ch23:central}. For a satellite at 1200\,km, find the slant range, central angle and
nadir angle at 20$^\circ$ elevation, and the fraction of the Earth covered above 20$^\circ$.
\item Compute the azimuth, elevation and range from your own location to a GEO satellite at your
longitude plus 30$^\circ$, and the polarisation skew angle you would need (look up or derive its
formula).
\item Derive the peak Doppler shift and Doppler rate~\eqref{eq:ch23:dopplermax} for an overhead
LEO pass. Evaluate them for Iridium (780\,km, 1.62\,GHz) and for a 350\,km very-low-orbit satellite at
2\,GHz. By roughly how much does the Earth's rotation change the answer for a station at the equator?
\item A 1\,m dish with 60\% efficiency, 0.3\,dB feed loss and a 1\,dB-noise-figure LNB views a clear
sky of 25\,K at 12\,GHz. Compute its $G/T$. By how much does a 4\,dB rain fade reduce the downlink
$C/N$, including the rise in sky noise?
\item A bent-pipe link has $(C/N)_u=18$\,dB, $(C/N)_d=14$\,dB, $C/I_{\text{ASI}}=22$\,dB and
$C/I_{\text{IM}}=19$\,dB. Find the total $C/(N+I)$. Which single improvement of 3\,dB helps most?
If the payload were regenerative and each link used a code with a steep waterfall, how would the
answer change qualitatively?
\item Repeat the Ka-band ACM worked example for a site with $R_{0.01}=90$\,mm/h and a 45\,cm dish.
What availability can QPSK 1/4 achieve? What would DVB-S2X's very-low-SNR MODCODs buy?
\item A transponder of 36\,MHz and saturated EIRP 50\,dBW is shared by SCPC carriers of
2\,Mbaud each, with 20\% guard bands, requiring $C/N=7$\,dB at receivers with $G/T=20$\,dB/K,
over a 205\,dB path. The multicarrier transponder must run at 4\,dB output back-off. Is it
power-limited or bandwidth-limited, and how many carriers fit?
\item Show that for pure ALOHA the throughput is $S=Ge^{-2G}$ and find its maximum. For a VSAT
network in which each of $10\,000$ terminals sends one 1000-bit transaction per minute on
a 256\,kb/s random-access carrier, what is $G$ and what fraction of packets collide?
\item Estimate the number of HARQ processes a stop-and-wait HARQ scheme would need to keep a GEO
NTN link busy with 1\,ms slots. Explain why Release~17 allows HARQ feedback to be disabled and what
replaces it.
\item Estimate the maximum rain attenuation at 40\,GHz exceeded 0.1\% of the time for a gateway with
$R_{0.01}=30$\,mm/h at 30$^\circ$ elevation. How far apart should two site-diversity gateways be,
and why does this help?
\item For the Voyager worked example, compute the data rate if Voyager used Ka-band (32\,GHz) with the
same antenna and power, a 34\,m antenna of 60\% efficiency at the ground and $T_{\text{sys}}=35$\,K,
ignoring weather. What limits this in practice?
\item \emph{(Simulation)} Use \texttt{satellite.leo\_pass} to generate the Doppler of a 550\,km pass at
2\,GHz, apply it to an NR-like OFDM signal with 15\,kHz subcarriers, and measure the
inter-carrier interference with and without pre-compensation from an ephemeris that is wrong by
1\,km in position.
\item \emph{(Simulation)} Extend the TWTA experiment of Figure~\ref{fig:ch23_apsk_vs_qam} to 32APSK
and 32QAM (cross). Add a simple symbol-level predistorter that adjusts each ring's radius and phase,
and find the new optimum total degradation.
\item \emph{(Simulation)} Using \texttt{rain\_atten\_p618}, generate a synthetic year of rain
attenuation for a Ka-band link (hint: draw from the exceedance distribution) and compute the
annual data volume delivered with CCM at 99.5\% availability and with ACM.
\end{enumerate}

\furtherreading
\begin{itemize}
\item G.~Maral, M.~Bousquet and Z.~Sun, \emph{Satellite Communications Systems: Systems, Techniques
and Technology}, 6th ed., Wiley, 2020. The comprehensive reference: orbits, link budgets, payloads,
access and networks.
\item T.~Pratt and J.~E.~Allnutt, \emph{Satellite Communications}, 3rd ed., Wiley, 2020. A clear
engineering textbook with many worked link budgets.
\item L.~J.~Ippolito, \emph{Satellite Communications Systems Engineering: Atmospheric Effects,
Satellite Link Design and System Performance}, 2nd ed., Wiley, 2017. The best guide to propagation
and the ITU-R models.
\item A.~C.~Clarke, ``Extra-Terrestrial Relays,'' \emph{Wireless World}, October 1945. Short and still
worth reading.
\item Recommendations ITU-R P.618 (Earth--space propagation), P.837 (rain rate), P.838 (specific
attenuation), P.839 (rain height), P.676 (gases) and S.580/S.465 (earth-station antenna patterns).
Free from the ITU; the source of the methods in Section~\ref{sec:ch23:linkbudget}.
\item ETSI EN~302~307-1 and -2, the DVB-S2 and DVB-S2X standards; ETSI EN~301~545-2, DVB-RCS2.
Remarkably readable standards.
\item A.~Morello and V.~Mignone, ``DVB-S2: The second generation standard for satellite broad-band
services,'' \emph{Proc. IEEE}, 94(1), 2006. The designers' own account of the choices.
\item A.~A.~M.~Saleh, ``Frequency-independent and frequency-dependent nonlinear models of TWT
amplifiers,'' \emph{IEEE Trans. Communications}, 29(11), 1981. The source of Equation~\eqref{eq:ch23:saleh}.
\item R.~De~Gaudenzi, A.~Guill\'en~i~F\`abregas and A.~Martinez, ``Turbo-coded APSK modulations design
for satellite broadband communications,'' \emph{Int. J. Satellite Communications and Networking},
2006. The analysis behind DVB-S2's APSK.
\item 3GPP TR~38.811 and TR~38.821, studies of NR for non-terrestrial networks (channel models,
architectures and solutions); TS~38.300 for the NTN overview in the specifications.
\item J.~H.~Yuen (ed.), \emph{Deep Space Telecommunications Systems Engineering}, JPL, 1982, and the
JPL DESCANSO monograph series (free online). How deep-space links are designed.
\end{itemize}
